% !TEX program = pdflatex
\documentclass[11pt,a4paper]{article}

% --- Packages ---
\usepackage[utf8]{inputenc}
\usepackage[T2A,T1]{fontenc}
\usepackage[ukrainian,english]{babel}
\usepackage{amsmath,amssymb}
\usepackage{graphicx}
\usepackage{booktabs}
\usepackage{multirow}
\usepackage{hyperref}
\usepackage{xcolor}
\usepackage{natbib}
\usepackage{url}
\usepackage{microtype}
\usepackage{enumitem}
\usepackage{caption}
\usepackage{subcaption}
\usepackage[margin=1in]{geometry}
\usepackage{float}
\usepackage{array}
\usepackage{tabularx}
\usepackage{pgfplots}
\usepackage{pgfplotstable}
\pgfplotsset{compat=1.18}
\usepgfplotslibrary{colorbrewer}

\hypersetup{
    colorlinks=true,
    linkcolor=blue!70!black,
    citecolor=blue!70!black,
    urlcolor=blue!70!black,
}

\newcolumntype{R}[1]{>{\raggedleft\arraybackslash}p{#1}}

% --- Title ---
\title{The Tokenizer Tax Across 24 European Languages:\\
Domain Invariance, Cross-Lingual Few-Shot Effects,\\
and the Ukrainian Penalty}

\author{
  Volodymyr Ovcharov\thanks{Corresponding author: \texttt{volodymyr@legal.org.ua}} \\
  LEX AI Platform, legal.org.ua \\
  Kyiv, Ukraine
}

\date{May 2026}

\begin{document}

\maketitle

% ============================================================
% ABSTRACT
% ============================================================
\begin{abstract}
Tokenizer fertility -- the number of tokens per word -- varies across languages and imposes a hidden cost on non-English NLP, but how large is this ``tokenizer tax,'' and does it affect downstream performance? We measure fertility for seven foundation models across 24~European languages on identical EU legislative text, producing the first comprehensive tokenizer tax map for the continent. The tax follows a clear hierarchy: English (1.2 tokens/word), Romance (1.6--1.9), Germanic (2.0--2.4), Slavic (2.6--3.1), Uralic/Baltic (3.3--3.6), with Greek at the extreme (5.7). Ukrainian (3.75) pays 20--40\% more than cognate Slavic languages on identical content, suggesting that its penalty combines morphological complexity with pre-training underrepresentation. We then show that fertility rankings are domain-invariant: the 1.6$\times$ spread between the most and least efficient tokenizers is identical on legal and news text. To test whether the tokenizer tax extends to downstream task quality, we evaluate few-shot classification on SIB-200 across four Slavic languages. The few-shot effect turns out to be model-intrinsic, not language-dependent: the same models that benefit from demonstrations in Ukrainian also benefit in Polish, Russian, and Czech, confirming recent mechanistic findings \citep{ovcharov2026fewshot}. A linguistic competence evaluation on the ULP benchmark shows fertility weakly predicts grammatical accuracy, but architecture matters more. Total experiment cost: \$4.05.
\end{abstract}

\noindent\textbf{Keywords:} tokenizer fertility, language tax, cross-lingual evaluation, few-shot prompting, morphologically rich languages, Ukrainian NLP, Slavic languages

% ============================================================
% 1. INTRODUCTION
% ============================================================
\section{Introduction}

Every token costs money. When a tokenizer fragments Ukrainian text into 3.75 tokens per word but handles English at 1.24, the same API call costs three times more -- a hidden ``tokenizer tax'' that penalizes billions of speakers of morphologically rich languages. \citet{petrov2024language} and \citet{ahia2023all} documented this disparity, but existing measurements cover individual languages on heterogeneous text, making it impossible to disentangle language effects from content effects.

This paper provides three things the literature lacks. First, a \textbf{controlled cross-lingual fertility map}: we measure seven foundation models on identical EU legislative text in 24~European languages, isolating the tokenizer's contribution from content variation. Second, a \textbf{domain invariance test}: we compare fertility on legal and news text to determine whether a single measurement suffices across applications. Third, a \textbf{downstream consequence test}: we evaluate few-shot classification across four Slavic languages to determine whether the tokenizer tax affects task performance.

\citet{ovcharov2026tokenizer} showed that tokenizer fertility varies by 1.6$\times$ across models on Ukrainian legal text and that few-shot prompting can degrade performance by up to 26~pp. \citet{ovcharov2026fewshot} subsequently explained the mechanistic basis of this degradation using representation analysis on 12~open-weight models. We extend these findings along the cross-lingual axis:

\begin{enumerate}[leftmargin=*]
    \item \textbf{Domain invariance} (Experiment~1): The 1.6$\times$ fertility spread between models is identical on legal and news text. A single measurement predicts cost across all domains.

    \item \textbf{Cross-lingual fertility map} (Experiment~2): Fertility across 24~EU languages on identical legal text, revealing a 4.6$\times$ spread from English (1.24) to Greek (5.66). Ukrainian pays 20--40\% more than cognate Slavic languages.

    \item \textbf{Cross-lingual few-shot} (Experiment~3): Few-shot classification on SIB-200 across Ukrainian, Polish, Russian, and Czech. The effect is model-intrinsic: the same models that benefit in one language benefit in all four, confirming the mechanistic findings of \citet{ovcharov2026fewshot}.

    \item \textbf{Linguistic competence} (Experiment~4): ULP benchmark (347 grammar questions) tests whether fertility predicts grammatical accuracy.
\end{enumerate}

% ============================================================
% 2. RELATED WORK
% ============================================================
\section{Related Work}

\subsection{Few-Shot Learning and Its Limits}

\citet{brown2020language} established few-shot in-context learning as a core capability of large language models. Subsequent work has shown that few-shot performance depends on example selection \citep{liu2022makes}, format consistency \citep{min2022rethinking}, and label distribution \citep{zhao2021calibrate}. \citet{min2022rethinking} demonstrated that even random labels can improve performance, suggesting that demonstrations primarily specify task format rather than input--output mappings.

However, these studies focus overwhelmingly on English. \citet{lai2023chatgpt} found significant performance variation across languages but did not systematically investigate few-shot effects. \citet{ovcharov2026tokenizer} documented systematic few-shot degradation on Ukrainian legal text. \citet{ovcharov2026fewshot} explained this mechanistically using representation analysis on 12~open-weight models, introducing a \emph{content delta} metric that predicts few-shot benefit ($\rho = +0.65$). Our work extends the cross-lingual dimension: testing whether the same patterns hold across four Slavic languages and 24~European languages.

\subsection{Tokenizer Fertility and Multilingual Fairness}

\citet{petrov2024language} formalized the ``language tax'' imposed by suboptimal tokenization, showing 2--15$\times$ token cost variation across languages. \citet{ahia2023all} demonstrated that API cost varies by an order of magnitude across languages for equivalent content. \citet{rust2021good} showed that monolingual performance correlates with pre-training data proportion, with fertility as a proxy.

These studies examine general-domain text or single languages. \citet{niklaus2024lextreme} introduced LEXTREME, a multi-lingual legal benchmark, but did not measure tokenizer fertility or few-shot effects. We provide the first controlled cross-lingual comparison: identical EU legal text in 24~languages, isolating the tokenizer's contribution from content variation. We also test domain invariance by comparing legal and news fertility for the same models.

\subsection{Morphologically Rich Languages in NLP}

Ukrainian, Polish, Russian, and Czech are Slavic languages with rich inflectional morphology: 7~cases, grammatical gender, and extensive verb conjugation. This morphological complexity interacts with subword tokenizers, potentially producing more fragmented representations that interfere with in-context pattern matching. \citet{conneau2020unsupervised} showed that multilingual model performance correlates with pre-training data volume, with low-resource languages suffering disproportionately. \citet{chaplynskyi2023ukrbruk} showed consistent underperformance of multilingual models on Ukrainian compared to English, a finding our cross-lingual experiments extend to the few-shot setting.

% ============================================================
% 3. METHODOLOGY
% ============================================================
\section{Methodology}

\subsection{Models}

We use the same seven models as \citet{ovcharov2026tokenizer}, all accessed via the AWS Bedrock API (Table~\ref{tab:models}). This ensures direct comparability with the legal text results.

\begin{table}[t]
\centering
\caption{Models evaluated. All accessed via AWS Bedrock in April--May 2026. MoE = mixture of experts; active = parameters per forward pass.}
\label{tab:models}
\begin{tabular}{llll}
\toprule
\textbf{Model} & \textbf{Provider} & \textbf{Architecture} & \textbf{Tokenizer} \\
\midrule
Llama 4 Maverick   & Meta      & 400B, 17B active, MoE & SentencePiece \\
Llama 3.3 70B      & Meta      & 70B dense              & SentencePiece \\
Mistral Large 3    & Mistral   & 675B, 41B active, MoE  & SentencePiece \\
Nemotron Super 3   & NVIDIA    & 120B, 12B active, hybrid & Custom \\
Nova Pro           & Amazon    & Undisclosed            & Proprietary \\
Qwen3 235B         & Qwen      & 235B, 22B active, MoE  & tiktoken-derived \\
Qwen3 32B          & Qwen      & 32B dense              & tiktoken-derived \\
\bottomrule
\end{tabular}
\end{table}

\subsection{Datasets}

\paragraph{SIB-200} \citep{adelani2024sib200} is a topic classification benchmark covering 205~languages with 1,000 examples each, annotated into 7~categories: science/technology, travel, politics, sports, health, entertainment, and geography. Critically, examples are parallel across languages (same \texttt{index\_id}), enabling paired cross-lingual comparisons. We use Ukrainian (\texttt{ukr\_Cyrl}), Polish (\texttt{pol\_Latn}), Russian (\texttt{rus\_Cyrl}), and Czech (\texttt{ces\_Latn}).

\paragraph{EU Acts in Ukrainian} \citep{francophonicEU} provides aligned EU legislative texts in 24~EU languages paired with Ukrainian, totaling over 3~million translation units. We sample 500~aligned segments per language pair and concatenate them into ${\sim}$6,000-character blocks for fertility measurement, matching the protocol of \citet{ovcharov2026tokenizer}.

\paragraph{ULP} \citep{galeshchuk2024ulp} is an expert-curated benchmark of 347~multiple-choice questions testing Ukrainian grammar and orthography, validated by professional linguists with 10--30~years of experience.

\subsection{Evaluation Protocol}

All evaluations use the AWS Bedrock \texttt{invoke\_model} API with temperature~0 for deterministic outputs, maintaining exact consistency with \citet{ovcharov2026tokenizer}. Provider-specific message formatting is preserved: Meta models use the Llama~3/4 prompt template, Amazon Nova uses the \texttt{messages-v1} schema, and remaining models (Qwen, Mistral, NVIDIA) use the standard messages format.

\paragraph{Fertility measurement.} We report the average ratio of API-reported input tokens to whitespace-delimited words. Because SIB-200 texts are short (typically 15--30 words), measuring fertility on individual sentences would be dominated by system prompt overhead. Following \citet{ovcharov2026tokenizer}, we concatenate texts into blocks of approximately 6,000 characters before measurement, yielding blocks of ${\sim}$840 words on average. For EU Acts, we concatenate aligned segments from each language into blocks of similar size. A trivial prompt (``Repeat the first word'') ensures the measurement captures tokenizer behavior on the input text rather than task-specific output.

\paragraph{Classification.} For SIB-200 topic classification, we report accuracy on the test split (204 parallel examples across all languages). Few-shot examples are drawn from the training split: one example per class (7~total for SIB-200). The prompt instructs the model to respond with only the English category name; responses are normalized via substring matching with a Ukrainian-to-English label map to handle models that respond in Ukrainian.

\paragraph{Linguistic competence.} For ULP, we report accuracy on 347 multiple-choice questions in zero-shot mode, and on 344 questions in few-shot mode (3~held out as demonstrations). The prompt presents each question with Ukrainian answer letters (\foreignlanguage{ukrainian}{А, Б, В, Г, Д}) and instructs the model to respond with only the letter.

% ============================================================
% 4. EXPERIMENTS AND RESULTS
% ============================================================
\section{Experiments and Results}

\subsection{Experiment 1: Cross-Domain Fertility}
\label{sec:exp1}

To test whether the fertility spread observed on legal text is domain-specific or a tokenizer property, we measured fertility on Ukrainian news text from SIB-200. We concatenated the 1,004 SIB-200 Ukrainian examples into 23 blocks of approximately 6,000 characters each, matching the measurement protocol of \citet{ovcharov2026tokenizer}.

Table~\ref{tab:crossdomain_fertility} presents the results alongside the legal text fertility values from \citet{ovcharov2026tokenizer}.

\begin{table}[t]
\centering
\caption{Cross-domain tokenizer fertility comparison: Ukrainian legal text \citep{ovcharov2026tokenizer} vs.\ Ukrainian news text (SIB-200). Fertility = tokens per whitespace-delimited word on ${\sim}$6,000-character blocks. Models sorted by news fertility.}
\label{tab:crossdomain_fertility}
\begin{tabular}{lR{1.4cm}R{1.4cm}R{1.4cm}}
\toprule
\textbf{Model} & \textbf{Legal} & \textbf{News} & \textbf{$\Delta$} \\
\midrule
Llama 4 Maverick   & 2.430 & 2.197 & $-$0.233 \\
Llama 3.3 70B      & 2.650 & 2.333 & $-$0.317 \\
Mistral Large 3    & 3.060 & 2.500 & $-$0.560 \\
Nemotron Super 3   & 3.080 & 2.516 & $-$0.564 \\
Nova Pro           & 2.850 & 3.269 & $+$0.419 \\
Qwen3 235B         & 3.894 & 3.578 & $-$0.316 \\
Qwen3 32B          & 3.902 & 3.583 & $-$0.319 \\
\midrule
\textit{Max/Min ratio} & \textit{1.61$\times$} & \textit{1.63$\times$} & \\
\bottomrule
\end{tabular}
\end{table}

Three patterns emerge. First, the \textbf{fertility spread is preserved}: the ratio between the least and most efficient tokenizers is 1.63$\times$ on news text (Llama~4 Maverick at 2.20 vs.\ Qwen3~32B at 3.58), virtually identical to the 1.61$\times$ spread on legal text. Second, \textbf{model rankings are stable}: both Qwen models remain the least efficient, both Llama models remain the most efficient, and the middle tier (Mistral, Nemotron) is preserved. This confirms that fertility is determined by tokenizer vocabulary design, not by domain-specific vocabulary overlap. Third, fertility is \textbf{consistently lower on news than legal text} for six of seven models ($\Delta$ ranges from $-$0.23 to $-$0.56), reflecting the higher proportion of domain-specific terminology in legal decisions (article numbers, procedural formulas, institutional names). The sole exception is Nova Pro ($+$0.42), whose proprietary tokenizer appears to handle legal boilerplate more efficiently than general news vocabulary.

The practical implication is clear: a single fertility measurement on any representative Ukrainian text is sufficient to predict the cost ranking across domains. Practitioners need not re-measure fertility for each new application.

\begin{figure}[t]
\centering
% Created by tikzDevice version 0.12.6 on 2026-05-15 12:02:42
% !TEX encoding = UTF-8 Unicode
\begin{tikzpicture}[x=1pt,y=1pt]
\definecolor{fillColor}{RGB}{255,255,255}
\path[use as bounding box,fill=fillColor,fill opacity=0.00] (0,0) rectangle (245.72,231.26);
\begin{scope}
\path[clip] (  3.81,  0.00) rectangle (241.91,231.26);
\definecolor{fillColor}{RGB}{255,255,255}

\path[fill=fillColor] (  3.81,  0.00) rectangle (241.91,231.26);
\end{scope}
\begin{scope}
\path[clip] ( 29.25, 22.61) rectangle (235.91,229.26);
\definecolor{drawColor}{gray}{0.92}

\path[draw=drawColor,line width= 0.5pt,line join=round] ( 29.25, 32.00) --
	(235.91, 32.00);

\path[draw=drawColor,line width= 0.5pt,line join=round] ( 29.25, 76.73) --
	(235.91, 76.73);

\path[draw=drawColor,line width= 0.5pt,line join=round] ( 29.25,121.46) --
	(235.91,121.46);

\path[draw=drawColor,line width= 0.5pt,line join=round] ( 29.25,166.19) --
	(235.91,166.19);

\path[draw=drawColor,line width= 0.5pt,line join=round] ( 29.25,210.92) --
	(235.91,210.92);

\path[draw=drawColor,line width= 0.5pt,line join=round] ( 38.65, 22.61) --
	( 38.65,229.26);

\path[draw=drawColor,line width= 0.5pt,line join=round] ( 83.38, 22.61) --
	( 83.38,229.26);

\path[draw=drawColor,line width= 0.5pt,line join=round] (128.11, 22.61) --
	(128.11,229.26);

\path[draw=drawColor,line width= 0.5pt,line join=round] (172.84, 22.61) --
	(172.84,229.26);

\path[draw=drawColor,line width= 0.5pt,line join=round] (217.57, 22.61) --
	(217.57,229.26);
\definecolor{drawColor}{gray}{0.60}

\path[draw=drawColor,line width= 0.5pt,dash pattern=on 4pt off 4pt ,line join=round] (-177.40,-184.05) -- (442.57,435.92);
\definecolor{drawColor}{RGB}{44,62,80}
\definecolor{fillColor}{RGB}{44,62,80}

\path[draw=drawColor,line width= 0.3pt,line join=round,line cap=round,fill=fillColor] ( 77.12, 49.62) circle (  2.97);

\path[draw=drawColor,line width= 0.3pt,line join=round,line cap=round,fill=fillColor] (208.80,173.62) circle (  2.97);

\path[draw=drawColor,line width= 0.3pt,line join=round,line cap=round,fill=fillColor] ( 96.80, 61.79) circle (  2.97);

\path[draw=drawColor,line width= 0.3pt,line join=round,line cap=round,fill=fillColor] (133.48, 76.73) circle (  2.97);

\path[draw=drawColor,line width= 0.3pt,line join=round,line cap=round,fill=fillColor] (208.09,173.17) circle (  2.97);

\path[draw=drawColor,line width= 0.3pt,line join=round,line cap=round,fill=fillColor] (135.27, 78.16) circle (  2.97);

\path[draw=drawColor,line width= 0.3pt,line join=round,line cap=round,fill=fillColor] (114.69,145.53) circle (  2.97);

\node[text=drawColor,anchor=base,inner sep=0pt, outer sep=0pt, scale=  0.74] at ( 77.12, 54.23) {Maverick};

\node[text=drawColor,anchor=base,inner sep=0pt, outer sep=0pt, scale=  0.74] at (208.80,178.23) {Qwen3 32B};

\node[text=drawColor,anchor=base,inner sep=0pt, outer sep=0pt, scale=  0.74] at ( 96.80, 66.40) {Llama 3.3};

\node[text=drawColor,anchor=base,inner sep=0pt, outer sep=0pt, scale=  0.74] at (133.48, 81.34) {Mistral};

\node[text=drawColor,anchor=base,inner sep=0pt, outer sep=0pt, scale=  0.74] at (208.09,177.78) {Qwen3 235B};

\node[text=drawColor,anchor=base,inner sep=0pt, outer sep=0pt, scale=  0.74] at (135.27, 82.77) {Nemotron};

\node[text=drawColor,anchor=base,inner sep=0pt, outer sep=0pt, scale=  0.74] at (114.69,150.14) {Nova Pro};
\end{scope}
\begin{scope}
\path[clip] (  0.00,  0.00) rectangle (245.72,231.26);
\definecolor{drawColor}{gray}{0.30}

\node[text=drawColor,anchor=base east,inner sep=0pt, outer sep=0pt, scale=  0.72] at ( 25.20, 29.52) {2.0};

\node[text=drawColor,anchor=base east,inner sep=0pt, outer sep=0pt, scale=  0.72] at ( 25.20, 74.25) {2.5};

\node[text=drawColor,anchor=base east,inner sep=0pt, outer sep=0pt, scale=  0.72] at ( 25.20,118.98) {3.0};

\node[text=drawColor,anchor=base east,inner sep=0pt, outer sep=0pt, scale=  0.72] at ( 25.20,163.71) {3.5};

\node[text=drawColor,anchor=base east,inner sep=0pt, outer sep=0pt, scale=  0.72] at ( 25.20,208.44) {4.0};
\end{scope}
\begin{scope}
\path[clip] (  0.00,  0.00) rectangle (245.72,231.26);
\definecolor{drawColor}{gray}{0.30}

\node[text=drawColor,anchor=base,inner sep=0pt, outer sep=0pt, scale=  0.72] at ( 38.65, 13.60) {2.0};

\node[text=drawColor,anchor=base,inner sep=0pt, outer sep=0pt, scale=  0.72] at ( 83.38, 13.60) {2.5};

\node[text=drawColor,anchor=base,inner sep=0pt, outer sep=0pt, scale=  0.72] at (128.11, 13.60) {3.0};

\node[text=drawColor,anchor=base,inner sep=0pt, outer sep=0pt, scale=  0.72] at (172.84, 13.60) {3.5};

\node[text=drawColor,anchor=base,inner sep=0pt, outer sep=0pt, scale=  0.72] at (217.57, 13.60) {4.0};
\end{scope}
\begin{scope}
\path[clip] (  0.00,  0.00) rectangle (245.72,231.26);
\definecolor{drawColor}{RGB}{0,0,0}

\node[text=drawColor,anchor=base,inner sep=0pt, outer sep=0pt, scale=  0.90] at (132.58,  3.75) {Fertility (legal text)};
\end{scope}
\begin{scope}
\path[clip] (  0.00,  0.00) rectangle (245.72,231.26);
\definecolor{drawColor}{RGB}{0,0,0}

\node[text=drawColor,rotate= 90.00,anchor=base,inner sep=0pt, outer sep=0pt, scale=  0.90] at ( 12.01,125.94) {Fertility (news text)};
\end{scope}
\end{tikzpicture}

\caption{Cross-domain fertility: legal text \citep{ovcharov2026tokenizer} vs.\ news text (SIB-200). The dashed line marks $y = x$; points below it indicate lower fertility on news. Model rankings are preserved (Spearman $\rho = 0.89$), confirming fertility is a tokenizer property. Nova~Pro is the sole outlier above the diagonal.}
\label{fig:crossdomain}
\end{figure}

\subsection{Experiment 2: Cross-Lingual Fertility}
\label{sec:exp2}

To contextualize Ukrainian's tokenizer penalty within the European language landscape, we measured fertility across all 24~EU languages using aligned EU legislation translations. Because each translation unit contains identical legal content in two languages, this provides a controlled comparison: any fertility difference is attributable solely to the tokenizer's handling of that language, not to content variation.

Table~\ref{tab:eu_fertility} presents fertility by language family for Qwen3~32B (the least efficient tokenizer). Figure~\ref{fig:eu_heatmap} shows the full 25-language $\times$ 7-model heatmap.

\begin{table}[t]
\centering
\caption{Tokenizer fertility on EU Acts legal text by language family (Qwen3~32B). Languages sorted by fertility within each family. English (1.24) serves as the baseline.}
\label{tab:eu_fertility}
\begin{tabular}{llR{1.8cm}}
\toprule
\textbf{Family} & \textbf{Language} & \textbf{Fertility} \\
\midrule
\multirow{2}{*}{English baseline} & English & 1.241 \\
                                   & Irish & 1.264 \\
\midrule
\multirow{4}{*}{Romance} & Spanish & 1.578 \\
                          & French & 1.622 \\
                          & Portuguese & 1.743 \\
                          & Italian & 1.857 \\
\midrule
\multirow{4}{*}{Germanic} & Dutch & 2.037 \\
                           & German & 2.182 \\
                           & Danish & 2.356 \\
                           & Swedish & 2.360 \\
\midrule
\multirow{5}{*}{Slavic} & Polish & 2.641 \\
                         & Slovenian & 2.651 \\
                         & Croatian & 2.684 \\
                         & Bulgarian & 2.766 \\
                         & Slovak & 3.065 \\
                         & Czech & 3.132 \\
\midrule
\multirow{3}{*}{Uralic} & Estonian & 3.351 \\
                         & Hungarian & 3.481 \\
                         & Finnish & 3.569 \\
\midrule
\multirow{2}{*}{Baltic} & Lithuanian & 3.267 \\
                         & Latvian & 3.435 \\
\midrule
Semitic & Maltese & 2.680 \\
\midrule
\textbf{Target} & \textbf{Ukrainian} & \textbf{3.753} \\
\midrule
Hellenic & Greek & 5.655 \\
\bottomrule
\end{tabular}
\end{table}

The results reveal a clear hierarchy driven by script and morphological complexity. Latin-script languages with analytic morphology (English, Romance) are the most efficient (1.2--1.9 tokens/word). Germanic languages cluster at 2.0--2.4. Slavic languages span 2.6--3.1, reflecting their rich inflectional morphology. Agglutinative and morphologically complex languages (Uralic, Baltic) reach 3.3--3.6.

Ukrainian (3.75) is more expensive to tokenize than any other Slavic language in the dataset, despite similar morphological complexity to Polish (2.64) or Czech (3.13). This gap---approximately 20--40\% higher than cognate Slavic languages---suggests that Ukrainian's disadvantage is not purely morphological but also reflects \textbf{underrepresentation in pre-training data}. Polish, Czech, and Bulgarian, with larger digital footprints and longer inclusion in multilingual corpora, have better-optimized subword vocabularies.

Greek (5.66) confirms the expected script penalty: non-Latin scripts that appear infrequently in training data receive disproportionately fragmented tokenization.

\begin{figure*}[t]
\centering
% Created by tikzDevice version 0.12.6 on 2026-05-15 12:04:26
% !TEX encoding = UTF-8 Unicode
\begin{tikzpicture}[x=1pt,y=1pt]
\definecolor{fillColor}{RGB}{255,255,255}
\path[use as bounding box,fill=fillColor,fill opacity=0.00] (0,0) rectangle (397.48,361.35);
\begin{scope}
\path[clip] (  0.00,  0.00) rectangle (397.48,361.35);
\definecolor{fillColor}{RGB}{255,255,255}

\path[fill=fillColor] ( -0.00,  0.00) rectangle (397.48,361.35);
\end{scope}
\begin{scope}
\path[clip] ( 39.39,  2.00) rectangle (354.87,328.55);
\definecolor{drawColor}{RGB}{255,255,255}
\definecolor{fillColor}{RGB}{255,245,161}

\path[draw=drawColor,line width= 0.3pt,fill=fillColor] ( 43.77,314.30) rectangle ( 87.59,327.26);
\definecolor{fillColor}{RGB}{255,209,100}

\path[draw=drawColor,line width= 0.3pt,fill=fillColor] ( 43.77, 16.25) rectangle ( 87.59, 29.21);
\definecolor{fillColor}{RGB}{255,173,76}

\path[draw=drawColor,line width= 0.3pt,fill=fillColor] (306.67,119.92) rectangle (350.49,132.88);
\definecolor{fillColor}{RGB}{254,153,66}

\path[draw=drawColor,line width= 0.3pt,fill=fillColor] (306.67, 94.01) rectangle (350.49,106.96);
\definecolor{fillColor}{RGB}{254,195,87}

\path[draw=drawColor,line width= 0.3pt,fill=fillColor] (306.67,197.67) rectangle (350.49,210.63);
\definecolor{fillColor}{RGB}{254,204,92}

\path[draw=drawColor,line width= 0.3pt,fill=fillColor] (306.67,223.59) rectangle (350.49,236.55);
\definecolor{fillColor}{RGB}{191,4,38}

\path[draw=drawColor,line width= 0.3pt,fill=fillColor] (306.67,  3.30) rectangle (350.49, 16.25);
\definecolor{fillColor}{RGB}{255,244,161}

\path[draw=drawColor,line width= 0.3pt,fill=fillColor] (306.67,314.30) rectangle (350.49,327.26);
\definecolor{fillColor}{RGB}{255,230,136}

\path[draw=drawColor,line width= 0.3pt,fill=fillColor] (306.67,288.38) rectangle (350.49,301.34);
\definecolor{fillColor}{RGB}{253,141,60}

\path[draw=drawColor,line width= 0.3pt,fill=fillColor] (306.67, 68.09) rectangle (350.49, 81.05);
\definecolor{fillColor}{RGB}{251,128,54}

\path[draw=drawColor,line width= 0.3pt,fill=fillColor] (306.67, 29.21) rectangle (350.49, 42.17);
\definecolor{fillColor}{RGB}{255,228,133}

\path[draw=drawColor,line width= 0.3pt,fill=fillColor] (306.67,275.42) rectangle (350.49,288.38);
\definecolor{fillColor}{RGB}{255,243,159}

\path[draw=drawColor,line width= 0.3pt,fill=fillColor] (306.67,301.34) rectangle (350.49,314.30);
\definecolor{fillColor}{RGB}{255,177,78}

\path[draw=drawColor,line width= 0.3pt,fill=fillColor] (306.67,132.88) rectangle (350.49,145.84);
\definecolor{fillColor}{RGB}{252,133,57}

\path[draw=drawColor,line width= 0.3pt,fill=fillColor] (306.67, 42.17) rectangle (350.49, 55.13);
\definecolor{fillColor}{RGB}{255,218,116}

\path[draw=drawColor,line width= 0.3pt,fill=fillColor] (306.67,249.51) rectangle (350.49,262.47);
\definecolor{fillColor}{RGB}{253,146,62}

\path[draw=drawColor,line width= 0.3pt,fill=fillColor] (306.67, 81.05) rectangle (350.49, 94.01);
\definecolor{fillColor}{RGB}{252,136,58}

\path[draw=drawColor,line width= 0.3pt,fill=fillColor] (306.67, 55.13) rectangle (350.49, 68.09);
\definecolor{fillColor}{RGB}{255,178,78}

\path[draw=drawColor,line width= 0.3pt,fill=fillColor] (306.67,145.84) rectangle (350.49,158.80);
\definecolor{fillColor}{RGB}{255,210,102}

\path[draw=drawColor,line width= 0.3pt,fill=fillColor] (306.67,236.55) rectangle (350.49,249.51);
\definecolor{fillColor}{RGB}{255,180,79}

\path[draw=drawColor,line width= 0.3pt,fill=fillColor] (306.67,171.76) rectangle (350.49,184.71);
\definecolor{fillColor}{RGB}{255,223,124}

\path[draw=drawColor,line width= 0.3pt,fill=fillColor] (306.67,262.47) rectangle (350.49,275.42);
\definecolor{fillColor}{RGB}{254,195,87}

\path[draw=drawColor,line width= 0.3pt,fill=fillColor] (306.67,210.63) rectangle (350.49,223.59);
\definecolor{fillColor}{RGB}{254,157,68}

\path[draw=drawColor,line width= 0.3pt,fill=fillColor] (306.67,106.96) rectangle (350.49,119.92);
\definecolor{fillColor}{RGB}{255,179,79}

\path[draw=drawColor,line width= 0.3pt,fill=fillColor] (306.67,158.80) rectangle (350.49,171.76);
\definecolor{fillColor}{RGB}{254,194,87}

\path[draw=drawColor,line width= 0.3pt,fill=fillColor] (306.67,184.71) rectangle (350.49,197.67);
\definecolor{fillColor}{RGB}{249,117,50}

\path[draw=drawColor,line width= 0.3pt,fill=fillColor] (306.67, 16.25) rectangle (350.49, 29.21);
\definecolor{fillColor}{RGB}{255,186,82}

\path[draw=drawColor,line width= 0.3pt,fill=fillColor] ( 87.59,119.92) rectangle (131.40,132.88);
\definecolor{fillColor}{RGB}{254,206,95}

\path[draw=drawColor,line width= 0.3pt,fill=fillColor] ( 87.59, 94.01) rectangle (131.40,106.96);
\definecolor{fillColor}{RGB}{254,197,88}

\path[draw=drawColor,line width= 0.3pt,fill=fillColor] ( 87.59,197.67) rectangle (131.40,210.63);
\definecolor{fillColor}{RGB}{254,205,93}

\path[draw=drawColor,line width= 0.3pt,fill=fillColor] ( 87.59,223.59) rectangle (131.40,236.55);
\definecolor{fillColor}{RGB}{255,188,84}

\path[draw=drawColor,line width= 0.3pt,fill=fillColor] ( 87.59,  3.30) rectangle (131.40, 16.25);
\definecolor{fillColor}{RGB}{255,246,163}

\path[draw=drawColor,line width= 0.3pt,fill=fillColor] ( 87.59,314.30) rectangle (131.40,327.26);
\definecolor{fillColor}{RGB}{255,230,137}

\path[draw=drawColor,line width= 0.3pt,fill=fillColor] ( 87.59,288.38) rectangle (131.40,301.34);
\definecolor{fillColor}{RGB}{253,145,62}

\path[draw=drawColor,line width= 0.3pt,fill=fillColor] ( 87.59, 68.09) rectangle (131.40, 81.05);
\definecolor{fillColor}{RGB}{252,132,56}

\path[draw=drawColor,line width= 0.3pt,fill=fillColor] ( 87.59, 29.21) rectangle (131.40, 42.17);
\definecolor{fillColor}{RGB}{255,228,133}

\path[draw=drawColor,line width= 0.3pt,fill=fillColor] ( 87.59,275.42) rectangle (131.40,288.38);
\definecolor{fillColor}{RGB}{255,245,161}

\path[draw=drawColor,line width= 0.3pt,fill=fillColor] ( 87.59,301.34) rectangle (131.40,314.30);
\definecolor{fillColor}{RGB}{255,184,81}

\path[draw=drawColor,line width= 0.3pt,fill=fillColor] ( 87.59,132.88) rectangle (131.40,145.84);
\definecolor{fillColor}{RGB}{252,137,58}

\path[draw=drawColor,line width= 0.3pt,fill=fillColor] ( 87.59, 42.17) rectangle (131.40, 55.13);
\definecolor{fillColor}{RGB}{255,219,118}

\path[draw=drawColor,line width= 0.3pt,fill=fillColor] ( 87.59,249.51) rectangle (131.40,262.47);
\definecolor{fillColor}{RGB}{253,141,60}

\path[draw=drawColor,line width= 0.3pt,fill=fillColor] ( 87.59, 81.05) rectangle (131.40, 94.01);
\definecolor{fillColor}{RGB}{253,138,59}

\path[draw=drawColor,line width= 0.3pt,fill=fillColor] ( 87.59, 55.13) rectangle (131.40, 68.09);
\definecolor{fillColor}{RGB}{255,171,75}

\path[draw=drawColor,line width= 0.3pt,fill=fillColor] ( 87.59,145.84) rectangle (131.40,158.80);
\definecolor{fillColor}{RGB}{255,211,105}

\path[draw=drawColor,line width= 0.3pt,fill=fillColor] ( 87.59,236.55) rectangle (131.40,249.51);
\definecolor{fillColor}{RGB}{255,171,75}

\path[draw=drawColor,line width= 0.3pt,fill=fillColor] ( 87.59,171.76) rectangle (131.40,184.71);
\definecolor{fillColor}{RGB}{255,223,124}

\path[draw=drawColor,line width= 0.3pt,fill=fillColor] ( 87.59,262.47) rectangle (131.40,275.42);
\definecolor{fillColor}{RGB}{254,197,88}

\path[draw=drawColor,line width= 0.3pt,fill=fillColor] ( 87.59,210.63) rectangle (131.40,223.59);
\definecolor{fillColor}{RGB}{255,182,80}

\path[draw=drawColor,line width= 0.3pt,fill=fillColor] ( 87.59,106.96) rectangle (131.40,119.92);
\definecolor{fillColor}{RGB}{255,187,83}

\path[draw=drawColor,line width= 0.3pt,fill=fillColor] ( 87.59,158.80) rectangle (131.40,171.76);
\definecolor{fillColor}{RGB}{254,197,88}

\path[draw=drawColor,line width= 0.3pt,fill=fillColor] ( 87.59,184.71) rectangle (131.40,197.67);
\definecolor{fillColor}{RGB}{254,202,91}

\path[draw=drawColor,line width= 0.3pt,fill=fillColor] ( 87.59, 16.25) rectangle (131.40, 29.21);
\definecolor{fillColor}{RGB}{254,206,95}

\path[draw=drawColor,line width= 0.3pt,fill=fillColor] (131.40,119.92) rectangle (175.22,132.88);
\definecolor{fillColor}{RGB}{254,197,88}

\path[draw=drawColor,line width= 0.3pt,fill=fillColor] (131.40, 94.01) rectangle (175.22,106.96);
\definecolor{fillColor}{RGB}{254,200,90}

\path[draw=drawColor,line width= 0.3pt,fill=fillColor] (131.40,197.67) rectangle (175.22,210.63);
\definecolor{fillColor}{RGB}{255,221,121}

\path[draw=drawColor,line width= 0.3pt,fill=fillColor] (131.40,223.59) rectangle (175.22,236.55);
\definecolor{fillColor}{RGB}{255,190,84}

\path[draw=drawColor,line width= 0.3pt,fill=fillColor] (131.40,  3.30) rectangle (175.22, 16.25);
\definecolor{fillColor}{RGB}{255,244,161}

\path[draw=drawColor,line width= 0.3pt,fill=fillColor] (131.40,314.30) rectangle (175.22,327.26);
\definecolor{fillColor}{RGB}{255,237,147}

\path[draw=drawColor,line width= 0.3pt,fill=fillColor] (131.40,288.38) rectangle (175.22,301.34);
\definecolor{fillColor}{RGB}{254,167,72}

\path[draw=drawColor,line width= 0.3pt,fill=fillColor] (131.40, 68.09) rectangle (175.22, 81.05);
\definecolor{fillColor}{RGB}{254,162,70}

\path[draw=drawColor,line width= 0.3pt,fill=fillColor] (131.40, 29.21) rectangle (175.22, 42.17);
\definecolor{fillColor}{RGB}{255,238,149}

\path[draw=drawColor,line width= 0.3pt,fill=fillColor] (131.40,275.42) rectangle (175.22,288.38);
\definecolor{fillColor}{RGB}{255,244,160}

\path[draw=drawColor,line width= 0.3pt,fill=fillColor] (131.40,301.34) rectangle (175.22,314.30);
\definecolor{fillColor}{RGB}{254,203,91}

\path[draw=drawColor,line width= 0.3pt,fill=fillColor] (131.40,132.88) rectangle (175.22,145.84);
\definecolor{fillColor}{RGB}{255,181,80}

\path[draw=drawColor,line width= 0.3pt,fill=fillColor] (131.40, 42.17) rectangle (175.22, 55.13);
\definecolor{fillColor}{RGB}{255,228,132}

\path[draw=drawColor,line width= 0.3pt,fill=fillColor] (131.40,249.51) rectangle (175.22,262.47);
\definecolor{fillColor}{RGB}{254,168,73}

\path[draw=drawColor,line width= 0.3pt,fill=fillColor] (131.40, 81.05) rectangle (175.22, 94.01);
\definecolor{fillColor}{RGB}{254,165,71}

\path[draw=drawColor,line width= 0.3pt,fill=fillColor] (131.40, 55.13) rectangle (175.22, 68.09);
\definecolor{fillColor}{RGB}{255,171,75}

\path[draw=drawColor,line width= 0.3pt,fill=fillColor] (131.40,145.84) rectangle (175.22,158.80);
\definecolor{fillColor}{RGB}{255,218,116}

\path[draw=drawColor,line width= 0.3pt,fill=fillColor] (131.40,236.55) rectangle (175.22,249.51);
\definecolor{fillColor}{RGB}{254,192,86}

\path[draw=drawColor,line width= 0.3pt,fill=fillColor] (131.40,171.76) rectangle (175.22,184.71);
\definecolor{fillColor}{RGB}{255,230,137}

\path[draw=drawColor,line width= 0.3pt,fill=fillColor] (131.40,262.47) rectangle (175.22,275.42);
\definecolor{fillColor}{RGB}{255,210,102}

\path[draw=drawColor,line width= 0.3pt,fill=fillColor] (131.40,210.63) rectangle (175.22,223.59);
\definecolor{fillColor}{RGB}{255,189,84}

\path[draw=drawColor,line width= 0.3pt,fill=fillColor] (131.40,106.96) rectangle (175.22,119.92);
\definecolor{fillColor}{RGB}{254,200,90}

\path[draw=drawColor,line width= 0.3pt,fill=fillColor] (131.40,158.80) rectangle (175.22,171.76);
\definecolor{fillColor}{RGB}{254,204,92}

\path[draw=drawColor,line width= 0.3pt,fill=fillColor] (131.40,184.71) rectangle (175.22,197.67);
\definecolor{fillColor}{RGB}{255,184,81}

\path[draw=drawColor,line width= 0.3pt,fill=fillColor] (131.40, 16.25) rectangle (175.22, 29.21);
\definecolor{fillColor}{RGB}{255,218,115}

\path[draw=drawColor,line width= 0.3pt,fill=fillColor] ( 43.77,119.92) rectangle ( 87.59,132.88);
\definecolor{fillColor}{RGB}{255,209,100}

\path[draw=drawColor,line width= 0.3pt,fill=fillColor] ( 43.77, 94.01) rectangle ( 87.59,106.96);
\definecolor{fillColor}{RGB}{254,207,98}

\path[draw=drawColor,line width= 0.3pt,fill=fillColor] ( 43.77,197.67) rectangle ( 87.59,210.63);
\definecolor{fillColor}{RGB}{255,226,130}

\path[draw=drawColor,line width= 0.3pt,fill=fillColor] ( 43.77,223.59) rectangle ( 87.59,236.55);
\definecolor{fillColor}{RGB}{254,200,90}

\path[draw=drawColor,line width= 0.3pt,fill=fillColor] ( 43.77,  3.30) rectangle ( 87.59, 16.25);
\definecolor{fillColor}{RGB}{255,241,155}

\path[draw=drawColor,line width= 0.3pt,fill=fillColor] ( 43.77,288.38) rectangle ( 87.59,301.34);
\definecolor{fillColor}{RGB}{255,173,75}

\path[draw=drawColor,line width= 0.3pt,fill=fillColor] ( 43.77, 68.09) rectangle ( 87.59, 81.05);
\definecolor{fillColor}{RGB}{255,175,76}

\path[draw=drawColor,line width= 0.3pt,fill=fillColor] ( 43.77, 29.21) rectangle ( 87.59, 42.17);
\definecolor{fillColor}{RGB}{255,239,152}

\path[draw=drawColor,line width= 0.3pt,fill=fillColor] ( 43.77,275.42) rectangle ( 87.59,288.38);
\definecolor{fillColor}{RGB}{255,243,158}

\path[draw=drawColor,line width= 0.3pt,fill=fillColor] ( 43.77,301.34) rectangle ( 87.59,314.30);
\definecolor{fillColor}{RGB}{254,203,92}

\path[draw=drawColor,line width= 0.3pt,fill=fillColor] ( 43.77,132.88) rectangle ( 87.59,145.84);
\definecolor{fillColor}{RGB}{254,198,89}

\path[draw=drawColor,line width= 0.3pt,fill=fillColor] ( 43.77, 42.17) rectangle ( 87.59, 55.13);
\definecolor{fillColor}{RGB}{255,235,144}

\path[draw=drawColor,line width= 0.3pt,fill=fillColor] ( 43.77,249.51) rectangle ( 87.59,262.47);
\definecolor{fillColor}{RGB}{255,186,82}

\path[draw=drawColor,line width= 0.3pt,fill=fillColor] ( 43.77, 81.05) rectangle ( 87.59, 94.01);

\path[draw=drawColor,line width= 0.3pt,fill=fillColor] ( 43.77, 55.13) rectangle ( 87.59, 68.09);
\definecolor{fillColor}{RGB}{255,176,77}

\path[draw=drawColor,line width= 0.3pt,fill=fillColor] ( 43.77,145.84) rectangle ( 87.59,158.80);
\definecolor{fillColor}{RGB}{255,228,133}

\path[draw=drawColor,line width= 0.3pt,fill=fillColor] ( 43.77,236.55) rectangle ( 87.59,249.51);
\definecolor{fillColor}{RGB}{255,214,109}

\path[draw=drawColor,line width= 0.3pt,fill=fillColor] ( 43.77,171.76) rectangle ( 87.59,184.71);
\definecolor{fillColor}{RGB}{255,238,149}

\path[draw=drawColor,line width= 0.3pt,fill=fillColor] ( 43.77,262.47) rectangle ( 87.59,275.42);
\definecolor{fillColor}{RGB}{255,222,123}

\path[draw=drawColor,line width= 0.3pt,fill=fillColor] ( 43.77,210.63) rectangle ( 87.59,223.59);
\definecolor{fillColor}{RGB}{254,204,93}

\path[draw=drawColor,line width= 0.3pt,fill=fillColor] ( 43.77,106.96) rectangle ( 87.59,119.92);
\definecolor{fillColor}{RGB}{254,203,92}

\path[draw=drawColor,line width= 0.3pt,fill=fillColor] ( 43.77,158.80) rectangle ( 87.59,171.76);
\definecolor{fillColor}{RGB}{255,216,114}

\path[draw=drawColor,line width= 0.3pt,fill=fillColor] ( 43.77,184.71) rectangle ( 87.59,197.67);
\definecolor{fillColor}{RGB}{255,173,76}

\path[draw=drawColor,line width= 0.3pt,fill=fillColor] (262.86,119.92) rectangle (306.67,132.88);
\definecolor{fillColor}{RGB}{254,153,66}

\path[draw=drawColor,line width= 0.3pt,fill=fillColor] (262.86, 94.01) rectangle (306.67,106.96);
\definecolor{fillColor}{RGB}{254,195,87}

\path[draw=drawColor,line width= 0.3pt,fill=fillColor] (262.86,197.67) rectangle (306.67,210.63);
\definecolor{fillColor}{RGB}{254,204,92}

\path[draw=drawColor,line width= 0.3pt,fill=fillColor] (262.86,223.59) rectangle (306.67,236.55);
\definecolor{fillColor}{RGB}{192,5,38}

\path[draw=drawColor,line width= 0.3pt,fill=fillColor] (262.86,  3.30) rectangle (306.67, 16.25);
\definecolor{fillColor}{RGB}{255,245,161}

\path[draw=drawColor,line width= 0.3pt,fill=fillColor] (262.86,314.30) rectangle (306.67,327.26);
\definecolor{fillColor}{RGB}{255,230,136}

\path[draw=drawColor,line width= 0.3pt,fill=fillColor] (262.86,288.38) rectangle (306.67,301.34);
\definecolor{fillColor}{RGB}{253,142,60}

\path[draw=drawColor,line width= 0.3pt,fill=fillColor] (262.86, 68.09) rectangle (306.67, 81.05);
\definecolor{fillColor}{RGB}{251,128,55}

\path[draw=drawColor,line width= 0.3pt,fill=fillColor] (262.86, 29.21) rectangle (306.67, 42.17);
\definecolor{fillColor}{RGB}{255,228,133}

\path[draw=drawColor,line width= 0.3pt,fill=fillColor] (262.86,275.42) rectangle (306.67,288.38);
\definecolor{fillColor}{RGB}{255,244,160}

\path[draw=drawColor,line width= 0.3pt,fill=fillColor] (262.86,301.34) rectangle (306.67,314.30);
\definecolor{fillColor}{RGB}{255,178,78}

\path[draw=drawColor,line width= 0.3pt,fill=fillColor] (262.86,132.88) rectangle (306.67,145.84);
\definecolor{fillColor}{RGB}{252,134,57}

\path[draw=drawColor,line width= 0.3pt,fill=fillColor] (262.86, 42.17) rectangle (306.67, 55.13);
\definecolor{fillColor}{RGB}{255,218,116}

\path[draw=drawColor,line width= 0.3pt,fill=fillColor] (262.86,249.51) rectangle (306.67,262.47);
\definecolor{fillColor}{RGB}{253,146,62}

\path[draw=drawColor,line width= 0.3pt,fill=fillColor] (262.86, 81.05) rectangle (306.67, 94.01);
\definecolor{fillColor}{RGB}{252,136,58}

\path[draw=drawColor,line width= 0.3pt,fill=fillColor] (262.86, 55.13) rectangle (306.67, 68.09);
\definecolor{fillColor}{RGB}{255,178,78}

\path[draw=drawColor,line width= 0.3pt,fill=fillColor] (262.86,145.84) rectangle (306.67,158.80);
\definecolor{fillColor}{RGB}{255,210,103}

\path[draw=drawColor,line width= 0.3pt,fill=fillColor] (262.86,236.55) rectangle (306.67,249.51);
\definecolor{fillColor}{RGB}{255,180,79}

\path[draw=drawColor,line width= 0.3pt,fill=fillColor] (262.86,171.76) rectangle (306.67,184.71);
\definecolor{fillColor}{RGB}{255,223,124}

\path[draw=drawColor,line width= 0.3pt,fill=fillColor] (262.86,262.47) rectangle (306.67,275.42);
\definecolor{fillColor}{RGB}{254,195,87}

\path[draw=drawColor,line width= 0.3pt,fill=fillColor] (262.86,210.63) rectangle (306.67,223.59);
\definecolor{fillColor}{RGB}{254,157,68}

\path[draw=drawColor,line width= 0.3pt,fill=fillColor] (262.86,106.96) rectangle (306.67,119.92);
\definecolor{fillColor}{RGB}{255,179,79}

\path[draw=drawColor,line width= 0.3pt,fill=fillColor] (262.86,158.80) rectangle (306.67,171.76);
\definecolor{fillColor}{RGB}{254,195,87}

\path[draw=drawColor,line width= 0.3pt,fill=fillColor] (262.86,184.71) rectangle (306.67,197.67);
\definecolor{fillColor}{RGB}{249,117,50}

\path[draw=drawColor,line width= 0.3pt,fill=fillColor] (262.86, 16.25) rectangle (306.67, 29.21);
\definecolor{fillColor}{RGB}{254,205,94}

\path[draw=drawColor,line width= 0.3pt,fill=fillColor] (175.22,119.92) rectangle (219.04,132.88);
\definecolor{fillColor}{RGB}{254,196,88}

\path[draw=drawColor,line width= 0.3pt,fill=fillColor] (175.22, 94.01) rectangle (219.04,106.96);
\definecolor{fillColor}{RGB}{254,199,89}

\path[draw=drawColor,line width= 0.3pt,fill=fillColor] (175.22,197.67) rectangle (219.04,210.63);
\definecolor{fillColor}{RGB}{255,220,120}

\path[draw=drawColor,line width= 0.3pt,fill=fillColor] (175.22,223.59) rectangle (219.04,236.55);
\definecolor{fillColor}{RGB}{255,187,83}

\path[draw=drawColor,line width= 0.3pt,fill=fillColor] (175.22,  3.30) rectangle (219.04, 16.25);
\definecolor{fillColor}{RGB}{255,244,160}

\path[draw=drawColor,line width= 0.3pt,fill=fillColor] (175.22,314.30) rectangle (219.04,327.26);
\definecolor{fillColor}{RGB}{255,236,146}

\path[draw=drawColor,line width= 0.3pt,fill=fillColor] (175.22,288.38) rectangle (219.04,301.34);
\definecolor{fillColor}{RGB}{254,165,71}

\path[draw=drawColor,line width= 0.3pt,fill=fillColor] (175.22, 68.09) rectangle (219.04, 81.05);
\definecolor{fillColor}{RGB}{254,160,69}

\path[draw=drawColor,line width= 0.3pt,fill=fillColor] (175.22, 29.21) rectangle (219.04, 42.17);
\definecolor{fillColor}{RGB}{255,237,148}

\path[draw=drawColor,line width= 0.3pt,fill=fillColor] (175.22,275.42) rectangle (219.04,288.38);
\definecolor{fillColor}{RGB}{255,243,158}

\path[draw=drawColor,line width= 0.3pt,fill=fillColor] (175.22,301.34) rectangle (219.04,314.30);
\definecolor{fillColor}{RGB}{254,202,91}

\path[draw=drawColor,line width= 0.3pt,fill=fillColor] (175.22,132.88) rectangle (219.04,145.84);
\definecolor{fillColor}{RGB}{255,180,79}

\path[draw=drawColor,line width= 0.3pt,fill=fillColor] (175.22, 42.17) rectangle (219.04, 55.13);
\definecolor{fillColor}{RGB}{255,227,131}

\path[draw=drawColor,line width= 0.3pt,fill=fillColor] (175.22,249.51) rectangle (219.04,262.47);
\definecolor{fillColor}{RGB}{254,167,73}

\path[draw=drawColor,line width= 0.3pt,fill=fillColor] (175.22, 81.05) rectangle (219.04, 94.01);
\definecolor{fillColor}{RGB}{254,164,71}

\path[draw=drawColor,line width= 0.3pt,fill=fillColor] (175.22, 55.13) rectangle (219.04, 68.09);
\definecolor{fillColor}{RGB}{255,170,74}

\path[draw=drawColor,line width= 0.3pt,fill=fillColor] (175.22,145.84) rectangle (219.04,158.80);
\definecolor{fillColor}{RGB}{255,217,115}

\path[draw=drawColor,line width= 0.3pt,fill=fillColor] (175.22,236.55) rectangle (219.04,249.51);
\definecolor{fillColor}{RGB}{255,191,85}

\path[draw=drawColor,line width= 0.3pt,fill=fillColor] (175.22,171.76) rectangle (219.04,184.71);
\definecolor{fillColor}{RGB}{255,229,136}

\path[draw=drawColor,line width= 0.3pt,fill=fillColor] (175.22,262.47) rectangle (219.04,275.42);
\definecolor{fillColor}{RGB}{255,209,101}

\path[draw=drawColor,line width= 0.3pt,fill=fillColor] (175.22,210.63) rectangle (219.04,223.59);
\definecolor{fillColor}{RGB}{255,188,83}

\path[draw=drawColor,line width= 0.3pt,fill=fillColor] (175.22,106.96) rectangle (219.04,119.92);
\definecolor{fillColor}{RGB}{254,200,90}

\path[draw=drawColor,line width= 0.3pt,fill=fillColor] (175.22,158.80) rectangle (219.04,171.76);
\definecolor{fillColor}{RGB}{254,203,91}

\path[draw=drawColor,line width= 0.3pt,fill=fillColor] (175.22,184.71) rectangle (219.04,197.67);
\definecolor{fillColor}{RGB}{255,183,81}

\path[draw=drawColor,line width= 0.3pt,fill=fillColor] (175.22, 16.25) rectangle (219.04, 29.21);
\definecolor{fillColor}{RGB}{254,170,74}

\path[draw=drawColor,line width= 0.3pt,fill=fillColor] (219.04,119.92) rectangle (262.86,132.88);
\definecolor{fillColor}{RGB}{254,165,71}

\path[draw=drawColor,line width= 0.3pt,fill=fillColor] (219.04, 94.01) rectangle (262.86,106.96);
\definecolor{fillColor}{RGB}{255,187,83}

\path[draw=drawColor,line width= 0.3pt,fill=fillColor] (219.04,197.67) rectangle (262.86,210.63);
\definecolor{fillColor}{RGB}{255,228,133}

\path[draw=drawColor,line width= 0.3pt,fill=fillColor] (219.04,223.59) rectangle (262.86,236.55);
\definecolor{fillColor}{RGB}{251,129,55}

\path[draw=drawColor,line width= 0.3pt,fill=fillColor] (219.04,  3.30) rectangle (262.86, 16.25);
\definecolor{fillColor}{RGB}{255,245,161}

\path[draw=drawColor,line width= 0.3pt,fill=fillColor] (219.04,314.30) rectangle (262.86,327.26);
\definecolor{fillColor}{RGB}{255,239,151}

\path[draw=drawColor,line width= 0.3pt,fill=fillColor] (219.04,288.38) rectangle (262.86,301.34);
\definecolor{fillColor}{RGB}{253,142,60}

\path[draw=drawColor,line width= 0.3pt,fill=fillColor] (219.04, 68.09) rectangle (262.86, 81.05);
\definecolor{fillColor}{RGB}{252,134,57}

\path[draw=drawColor,line width= 0.3pt,fill=fillColor] (219.04, 29.21) rectangle (262.86, 42.17);
\definecolor{fillColor}{RGB}{255,230,136}

\path[draw=drawColor,line width= 0.3pt,fill=fillColor] (219.04,275.42) rectangle (262.86,288.38);
\definecolor{fillColor}{RGB}{255,244,160}

\path[draw=drawColor,line width= 0.3pt,fill=fillColor] (219.04,301.34) rectangle (262.86,314.30);
\definecolor{fillColor}{RGB}{255,180,79}

\path[draw=drawColor,line width= 0.3pt,fill=fillColor] (219.04,132.88) rectangle (262.86,145.84);
\definecolor{fillColor}{RGB}{251,130,55}

\path[draw=drawColor,line width= 0.3pt,fill=fillColor] (219.04, 42.17) rectangle (262.86, 55.13);
\definecolor{fillColor}{RGB}{255,230,136}

\path[draw=drawColor,line width= 0.3pt,fill=fillColor] (219.04,249.51) rectangle (262.86,262.47);
\definecolor{fillColor}{RGB}{254,149,64}

\path[draw=drawColor,line width= 0.3pt,fill=fillColor] (219.04, 81.05) rectangle (262.86, 94.01);
\definecolor{fillColor}{RGB}{253,140,60}

\path[draw=drawColor,line width= 0.3pt,fill=fillColor] (219.04, 55.13) rectangle (262.86, 68.09);
\definecolor{fillColor}{RGB}{255,177,78}

\path[draw=drawColor,line width= 0.3pt,fill=fillColor] (219.04,145.84) rectangle (262.86,158.80);
\definecolor{fillColor}{RGB}{254,205,93}

\path[draw=drawColor,line width= 0.3pt,fill=fillColor] (219.04,236.55) rectangle (262.86,249.51);
\definecolor{fillColor}{RGB}{254,159,69}

\path[draw=drawColor,line width= 0.3pt,fill=fillColor] (219.04,171.76) rectangle (262.86,184.71);
\definecolor{fillColor}{RGB}{255,230,136}

\path[draw=drawColor,line width= 0.3pt,fill=fillColor] (219.04,262.47) rectangle (262.86,275.42);
\definecolor{fillColor}{RGB}{254,197,88}

\path[draw=drawColor,line width= 0.3pt,fill=fillColor] (219.04,210.63) rectangle (262.86,223.59);
\definecolor{fillColor}{RGB}{254,164,71}

\path[draw=drawColor,line width= 0.3pt,fill=fillColor] (219.04,106.96) rectangle (262.86,119.92);
\definecolor{fillColor}{RGB}{255,181,80}

\path[draw=drawColor,line width= 0.3pt,fill=fillColor] (219.04,158.80) rectangle (262.86,171.76);
\definecolor{fillColor}{RGB}{255,188,83}

\path[draw=drawColor,line width= 0.3pt,fill=fillColor] (219.04,184.71) rectangle (262.86,197.67);
\definecolor{fillColor}{RGB}{253,142,61}

\path[draw=drawColor,line width= 0.3pt,fill=fillColor] (219.04, 16.25) rectangle (262.86, 29.21);
\definecolor{drawColor}{gray}{0.20}

\node[text=drawColor,anchor=base,inner sep=0pt, outer sep=0pt, scale=  0.54] at ( 65.68,318.92) {1.2};

\node[text=drawColor,anchor=base,inner sep=0pt, outer sep=0pt, scale=  0.54] at ( 65.68, 20.87) {2.1};

\node[text=drawColor,anchor=base,inner sep=0pt, outer sep=0pt, scale=  0.54] at (328.58,124.54) {2.8};

\node[text=drawColor,anchor=base,inner sep=0pt, outer sep=0pt, scale=  0.54] at (328.58, 98.62) {3.1};

\node[text=drawColor,anchor=base,inner sep=0pt, outer sep=0pt, scale=  0.54] at (328.58,202.29) {2.4};

\node[text=drawColor,anchor=base,inner sep=0pt, outer sep=0pt, scale=  0.54] at (328.58,228.21) {2.2};

\node[text=drawColor,anchor=base,inner sep=0pt, outer sep=0pt, scale=  0.54] at (328.58,  7.91) {5.7};

\node[text=drawColor,anchor=base,inner sep=0pt, outer sep=0pt, scale=  0.54] at (328.58,318.92) {1.2};

\node[text=drawColor,anchor=base,inner sep=0pt, outer sep=0pt, scale=  0.54] at (328.58,293.00) {1.6};

\node[text=drawColor,anchor=base,inner sep=0pt, outer sep=0pt, scale=  0.54] at (328.58, 72.71) {3.4};

\node[text=drawColor,anchor=base,inner sep=0pt, outer sep=0pt, scale=  0.54] at (328.58, 33.83) {3.6};

\node[text=drawColor,anchor=base,inner sep=0pt, outer sep=0pt, scale=  0.54] at (328.58,280.04) {1.6};

\node[text=drawColor,anchor=base,inner sep=0pt, outer sep=0pt, scale=  0.54] at (328.58,305.96) {1.3};

\node[text=drawColor,anchor=base,inner sep=0pt, outer sep=0pt, scale=  0.54] at (328.58,137.50) {2.7};

\node[text=drawColor,anchor=base,inner sep=0pt, outer sep=0pt, scale=  0.54] at (328.58, 46.79) {3.5};

\node[text=drawColor,anchor=base,inner sep=0pt, outer sep=0pt, scale=  0.54] at (328.58,254.12) {1.9};

\node[text=drawColor,anchor=base,inner sep=0pt, outer sep=0pt, scale=  0.54] at (328.58, 85.66) {3.3};

\node[text=drawColor,anchor=base,inner sep=0pt, outer sep=0pt, scale=  0.54] at (328.58, 59.75) {3.4};

\node[text=drawColor,anchor=base,inner sep=0pt, outer sep=0pt, scale=  0.54] at (328.58,150.46) {2.7};

\node[text=drawColor,anchor=base,inner sep=0pt, outer sep=0pt, scale=  0.54] at (328.58,241.17) {2.0};

\node[text=drawColor,anchor=base,inner sep=0pt, outer sep=0pt, scale=  0.54] at (328.58,176.37) {2.6};

\node[text=drawColor,anchor=base,inner sep=0pt, outer sep=0pt, scale=  0.54] at (328.58,267.08) {1.7};

\node[text=drawColor,anchor=base,inner sep=0pt, outer sep=0pt, scale=  0.54] at (328.58,215.25) {2.4};

\node[text=drawColor,anchor=base,inner sep=0pt, outer sep=0pt, scale=  0.54] at (328.58,111.58) {3.1};

\node[text=drawColor,anchor=base,inner sep=0pt, outer sep=0pt, scale=  0.54] at (328.58,163.42) {2.7};

\node[text=drawColor,anchor=base,inner sep=0pt, outer sep=0pt, scale=  0.54] at (328.58,189.33) {2.4};

\node[text=drawColor,anchor=base,inner sep=0pt, outer sep=0pt, scale=  0.54] at (328.58, 20.87) {3.8};

\node[text=drawColor,anchor=base,inner sep=0pt, outer sep=0pt, scale=  0.54] at (109.49,124.54) {2.5};

\node[text=drawColor,anchor=base,inner sep=0pt, outer sep=0pt, scale=  0.54] at (109.49, 98.62) {2.1};

\node[text=drawColor,anchor=base,inner sep=0pt, outer sep=0pt, scale=  0.54] at (109.49,202.29) {2.3};

\node[text=drawColor,anchor=base,inner sep=0pt, outer sep=0pt, scale=  0.54] at (109.49,228.21) {2.2};

\node[text=drawColor,anchor=base,inner sep=0pt, outer sep=0pt, scale=  0.54] at (109.49,  7.91) {2.5};

\node[text=drawColor,anchor=base,inner sep=0pt, outer sep=0pt, scale=  0.54] at (109.49,318.92) {1.2};

\node[text=drawColor,anchor=base,inner sep=0pt, outer sep=0pt, scale=  0.54] at (109.49,293.00) {1.6};

\node[text=drawColor,anchor=base,inner sep=0pt, outer sep=0pt, scale=  0.54] at (109.49, 72.71) {3.3};

\node[text=drawColor,anchor=base,inner sep=0pt, outer sep=0pt, scale=  0.54] at (109.49, 33.83) {3.5};

\node[text=drawColor,anchor=base,inner sep=0pt, outer sep=0pt, scale=  0.54] at (109.49,280.04) {1.6};

\node[text=drawColor,anchor=base,inner sep=0pt, outer sep=0pt, scale=  0.54] at (109.49,305.96) {1.2};

\node[text=drawColor,anchor=base,inner sep=0pt, outer sep=0pt, scale=  0.54] at (109.49,137.50) {2.6};

\node[text=drawColor,anchor=base,inner sep=0pt, outer sep=0pt, scale=  0.54] at (109.49, 46.79) {3.4};

\node[text=drawColor,anchor=base,inner sep=0pt, outer sep=0pt, scale=  0.54] at (109.49,254.12) {1.8};

\node[text=drawColor,anchor=base,inner sep=0pt, outer sep=0pt, scale=  0.54] at (109.49, 85.66) {3.3};

\node[text=drawColor,anchor=base,inner sep=0pt, outer sep=0pt, scale=  0.54] at (109.49, 59.75) {3.4};

\node[text=drawColor,anchor=base,inner sep=0pt, outer sep=0pt, scale=  0.54] at (109.49,150.46) {2.8};

\node[text=drawColor,anchor=base,inner sep=0pt, outer sep=0pt, scale=  0.54] at (109.49,241.17) {2.0};

\node[text=drawColor,anchor=base,inner sep=0pt, outer sep=0pt, scale=  0.54] at (109.49,176.37) {2.8};

\node[text=drawColor,anchor=base,inner sep=0pt, outer sep=0pt, scale=  0.54] at (109.49,267.08) {1.7};

\node[text=drawColor,anchor=base,inner sep=0pt, outer sep=0pt, scale=  0.54] at (109.49,215.25) {2.3};

\node[text=drawColor,anchor=base,inner sep=0pt, outer sep=0pt, scale=  0.54] at (109.49,111.58) {2.6};

\node[text=drawColor,anchor=base,inner sep=0pt, outer sep=0pt, scale=  0.54] at (109.49,163.42) {2.5};

\node[text=drawColor,anchor=base,inner sep=0pt, outer sep=0pt, scale=  0.54] at (109.49,189.33) {2.3};

\node[text=drawColor,anchor=base,inner sep=0pt, outer sep=0pt, scale=  0.54] at (109.49, 20.87) {2.2};

\node[text=drawColor,anchor=base,inner sep=0pt, outer sep=0pt, scale=  0.54] at (153.31,124.54) {2.1};

\node[text=drawColor,anchor=base,inner sep=0pt, outer sep=0pt, scale=  0.54] at (153.31, 98.62) {2.3};

\node[text=drawColor,anchor=base,inner sep=0pt, outer sep=0pt, scale=  0.54] at (153.31,202.29) {2.2};

\node[text=drawColor,anchor=base,inner sep=0pt, outer sep=0pt, scale=  0.54] at (153.31,228.21) {1.8};

\node[text=drawColor,anchor=base,inner sep=0pt, outer sep=0pt, scale=  0.54] at (153.31,  7.91) {2.4};

\node[text=drawColor,anchor=base,inner sep=0pt, outer sep=0pt, scale=  0.54] at (153.31,318.92) {1.2};

\node[text=drawColor,anchor=base,inner sep=0pt, outer sep=0pt, scale=  0.54] at (153.31,293.00) {1.4};

\node[text=drawColor,anchor=base,inner sep=0pt, outer sep=0pt, scale=  0.54] at (153.31, 72.71) {2.9};

\node[text=drawColor,anchor=base,inner sep=0pt, outer sep=0pt, scale=  0.54] at (153.31, 33.83) {3.0};

\node[text=drawColor,anchor=base,inner sep=0pt, outer sep=0pt, scale=  0.54] at (153.31,280.04) {1.4};

\node[text=drawColor,anchor=base,inner sep=0pt, outer sep=0pt, scale=  0.54] at (153.31,305.96) {1.3};

\node[text=drawColor,anchor=base,inner sep=0pt, outer sep=0pt, scale=  0.54] at (153.31,137.50) {2.2};

\node[text=drawColor,anchor=base,inner sep=0pt, outer sep=0pt, scale=  0.54] at (153.31, 46.79) {2.6};

\node[text=drawColor,anchor=base,inner sep=0pt, outer sep=0pt, scale=  0.54] at (153.31,254.12) {1.6};

\node[text=drawColor,anchor=base,inner sep=0pt, outer sep=0pt, scale=  0.54] at (153.31, 85.66) {2.9};

\node[text=drawColor,anchor=base,inner sep=0pt, outer sep=0pt, scale=  0.54] at (153.31, 59.75) {2.9};

\node[text=drawColor,anchor=base,inner sep=0pt, outer sep=0pt, scale=  0.54] at (153.31,150.46) {2.8};

\node[text=drawColor,anchor=base,inner sep=0pt, outer sep=0pt, scale=  0.54] at (153.31,241.17) {1.9};

\node[text=drawColor,anchor=base,inner sep=0pt, outer sep=0pt, scale=  0.54] at (153.31,176.37) {2.4};

\node[text=drawColor,anchor=base,inner sep=0pt, outer sep=0pt, scale=  0.54] at (153.31,267.08) {1.6};

\node[text=drawColor,anchor=base,inner sep=0pt, outer sep=0pt, scale=  0.54] at (153.31,215.25) {2.0};

\node[text=drawColor,anchor=base,inner sep=0pt, outer sep=0pt, scale=  0.54] at (153.31,111.58) {2.5};

\node[text=drawColor,anchor=base,inner sep=0pt, outer sep=0pt, scale=  0.54] at (153.31,163.42) {2.2};

\node[text=drawColor,anchor=base,inner sep=0pt, outer sep=0pt, scale=  0.54] at (153.31,189.33) {2.2};

\node[text=drawColor,anchor=base,inner sep=0pt, outer sep=0pt, scale=  0.54] at (153.31, 20.87) {2.6};

\node[text=drawColor,anchor=base,inner sep=0pt, outer sep=0pt, scale=  0.54] at ( 65.68,124.54) {1.9};

\node[text=drawColor,anchor=base,inner sep=0pt, outer sep=0pt, scale=  0.54] at ( 65.68, 98.62) {2.1};

\node[text=drawColor,anchor=base,inner sep=0pt, outer sep=0pt, scale=  0.54] at ( 65.68,202.29) {2.1};

\node[text=drawColor,anchor=base,inner sep=0pt, outer sep=0pt, scale=  0.54] at ( 65.68,228.21) {1.7};

\node[text=drawColor,anchor=base,inner sep=0pt, outer sep=0pt, scale=  0.54] at ( 65.68,  7.91) {2.3};

\node[text=drawColor,anchor=base,inner sep=0pt, outer sep=0pt, scale=  0.54] at ( 65.68,293.00) {1.3};

\node[text=drawColor,anchor=base,inner sep=0pt, outer sep=0pt, scale=  0.54] at ( 65.68, 72.71) {2.8};

\node[text=drawColor,anchor=base,inner sep=0pt, outer sep=0pt, scale=  0.54] at ( 65.68, 33.83) {2.7};

\node[text=drawColor,anchor=base,inner sep=0pt, outer sep=0pt, scale=  0.54] at ( 65.68,280.04) {1.4};

\node[text=drawColor,anchor=base,inner sep=0pt, outer sep=0pt, scale=  0.54] at ( 65.68,305.96) {1.3};

\node[text=drawColor,anchor=base,inner sep=0pt, outer sep=0pt, scale=  0.54] at ( 65.68,137.50) {2.2};

\node[text=drawColor,anchor=base,inner sep=0pt, outer sep=0pt, scale=  0.54] at ( 65.68, 46.79) {2.3};

\node[text=drawColor,anchor=base,inner sep=0pt, outer sep=0pt, scale=  0.54] at ( 65.68,254.12) {1.5};

\node[text=drawColor,anchor=base,inner sep=0pt, outer sep=0pt, scale=  0.54] at ( 65.68, 85.66) {2.5};

\node[text=drawColor,anchor=base,inner sep=0pt, outer sep=0pt, scale=  0.54] at ( 65.68, 59.75) {2.5};

\node[text=drawColor,anchor=base,inner sep=0pt, outer sep=0pt, scale=  0.54] at ( 65.68,150.46) {2.7};

\node[text=drawColor,anchor=base,inner sep=0pt, outer sep=0pt, scale=  0.54] at ( 65.68,241.17) {1.6};

\node[text=drawColor,anchor=base,inner sep=0pt, outer sep=0pt, scale=  0.54] at ( 65.68,176.37) {1.9};

\node[text=drawColor,anchor=base,inner sep=0pt, outer sep=0pt, scale=  0.54] at ( 65.68,267.08) {1.4};

\node[text=drawColor,anchor=base,inner sep=0pt, outer sep=0pt, scale=  0.54] at ( 65.68,215.25) {1.8};

\node[text=drawColor,anchor=base,inner sep=0pt, outer sep=0pt, scale=  0.54] at ( 65.68,111.58) {2.2};

\node[text=drawColor,anchor=base,inner sep=0pt, outer sep=0pt, scale=  0.54] at ( 65.68,163.42) {2.2};

\node[text=drawColor,anchor=base,inner sep=0pt, outer sep=0pt, scale=  0.54] at ( 65.68,189.33) {1.9};

\node[text=drawColor,anchor=base,inner sep=0pt, outer sep=0pt, scale=  0.54] at (284.77,124.54) {2.8};

\node[text=drawColor,anchor=base,inner sep=0pt, outer sep=0pt, scale=  0.54] at (284.77, 98.62) {3.1};

\node[text=drawColor,anchor=base,inner sep=0pt, outer sep=0pt, scale=  0.54] at (284.77,202.29) {2.4};

\node[text=drawColor,anchor=base,inner sep=0pt, outer sep=0pt, scale=  0.54] at (284.77,228.21) {2.2};

\node[text=drawColor,anchor=base,inner sep=0pt, outer sep=0pt, scale=  0.54] at (284.77,  7.91) {5.6};

\node[text=drawColor,anchor=base,inner sep=0pt, outer sep=0pt, scale=  0.54] at (284.77,318.92) {1.2};

\node[text=drawColor,anchor=base,inner sep=0pt, outer sep=0pt, scale=  0.54] at (284.77,293.00) {1.6};

\node[text=drawColor,anchor=base,inner sep=0pt, outer sep=0pt, scale=  0.54] at (284.77, 72.71) {3.3};

\node[text=drawColor,anchor=base,inner sep=0pt, outer sep=0pt, scale=  0.54] at (284.77, 33.83) {3.6};

\node[text=drawColor,anchor=base,inner sep=0pt, outer sep=0pt, scale=  0.54] at (284.77,280.04) {1.6};

\node[text=drawColor,anchor=base,inner sep=0pt, outer sep=0pt, scale=  0.54] at (284.77,305.96) {1.3};

\node[text=drawColor,anchor=base,inner sep=0pt, outer sep=0pt, scale=  0.54] at (284.77,137.50) {2.7};

\node[text=drawColor,anchor=base,inner sep=0pt, outer sep=0pt, scale=  0.54] at (284.77, 46.79) {3.5};

\node[text=drawColor,anchor=base,inner sep=0pt, outer sep=0pt, scale=  0.54] at (284.77,254.12) {1.9};

\node[text=drawColor,anchor=base,inner sep=0pt, outer sep=0pt, scale=  0.54] at (284.77, 85.66) {3.3};

\node[text=drawColor,anchor=base,inner sep=0pt, outer sep=0pt, scale=  0.54] at (284.77, 59.75) {3.4};

\node[text=drawColor,anchor=base,inner sep=0pt, outer sep=0pt, scale=  0.54] at (284.77,150.46) {2.7};

\node[text=drawColor,anchor=base,inner sep=0pt, outer sep=0pt, scale=  0.54] at (284.77,241.17) {2.0};

\node[text=drawColor,anchor=base,inner sep=0pt, outer sep=0pt, scale=  0.54] at (284.77,176.37) {2.6};

\node[text=drawColor,anchor=base,inner sep=0pt, outer sep=0pt, scale=  0.54] at (284.77,267.08) {1.7};

\node[text=drawColor,anchor=base,inner sep=0pt, outer sep=0pt, scale=  0.54] at (284.77,215.25) {2.3};

\node[text=drawColor,anchor=base,inner sep=0pt, outer sep=0pt, scale=  0.54] at (284.77,111.58) {3.1};

\node[text=drawColor,anchor=base,inner sep=0pt, outer sep=0pt, scale=  0.54] at (284.77,163.42) {2.6};

\node[text=drawColor,anchor=base,inner sep=0pt, outer sep=0pt, scale=  0.54] at (284.77,189.33) {2.4};

\node[text=drawColor,anchor=base,inner sep=0pt, outer sep=0pt, scale=  0.54] at (284.77, 20.87) {3.7};

\node[text=drawColor,anchor=base,inner sep=0pt, outer sep=0pt, scale=  0.54] at (197.13,124.54) {2.2};

\node[text=drawColor,anchor=base,inner sep=0pt, outer sep=0pt, scale=  0.54] at (197.13, 98.62) {2.3};

\node[text=drawColor,anchor=base,inner sep=0pt, outer sep=0pt, scale=  0.54] at (197.13,202.29) {2.3};

\node[text=drawColor,anchor=base,inner sep=0pt, outer sep=0pt, scale=  0.54] at (197.13,228.21) {1.8};

\node[text=drawColor,anchor=base,inner sep=0pt, outer sep=0pt, scale=  0.54] at (197.13,  7.91) {2.5};

\node[text=drawColor,anchor=base,inner sep=0pt, outer sep=0pt, scale=  0.54] at (197.13,318.92) {1.3};

\node[text=drawColor,anchor=base,inner sep=0pt, outer sep=0pt, scale=  0.54] at (197.13,293.00) {1.4};

\node[text=drawColor,anchor=base,inner sep=0pt, outer sep=0pt, scale=  0.54] at (197.13, 72.71) {2.9};

\node[text=drawColor,anchor=base,inner sep=0pt, outer sep=0pt, scale=  0.54] at (197.13, 33.83) {3.0};

\node[text=drawColor,anchor=base,inner sep=0pt, outer sep=0pt, scale=  0.54] at (197.13,280.04) {1.4};

\node[text=drawColor,anchor=base,inner sep=0pt, outer sep=0pt, scale=  0.54] at (197.13,305.96) {1.3};

\node[text=drawColor,anchor=base,inner sep=0pt, outer sep=0pt, scale=  0.54] at (197.13,137.50) {2.2};

\node[text=drawColor,anchor=base,inner sep=0pt, outer sep=0pt, scale=  0.54] at (197.13, 46.79) {2.6};

\node[text=drawColor,anchor=base,inner sep=0pt, outer sep=0pt, scale=  0.54] at (197.13,254.12) {1.6};

\node[text=drawColor,anchor=base,inner sep=0pt, outer sep=0pt, scale=  0.54] at (197.13, 85.66) {2.9};

\node[text=drawColor,anchor=base,inner sep=0pt, outer sep=0pt, scale=  0.54] at (197.13, 59.75) {2.9};

\node[text=drawColor,anchor=base,inner sep=0pt, outer sep=0pt, scale=  0.54] at (197.13,150.46) {2.8};

\node[text=drawColor,anchor=base,inner sep=0pt, outer sep=0pt, scale=  0.54] at (197.13,241.17) {1.9};

\node[text=drawColor,anchor=base,inner sep=0pt, outer sep=0pt, scale=  0.54] at (197.13,176.37) {2.4};

\node[text=drawColor,anchor=base,inner sep=0pt, outer sep=0pt, scale=  0.54] at (197.13,267.08) {1.6};

\node[text=drawColor,anchor=base,inner sep=0pt, outer sep=0pt, scale=  0.54] at (197.13,215.25) {2.1};

\node[text=drawColor,anchor=base,inner sep=0pt, outer sep=0pt, scale=  0.54] at (197.13,111.58) {2.5};

\node[text=drawColor,anchor=base,inner sep=0pt, outer sep=0pt, scale=  0.54] at (197.13,163.42) {2.3};

\node[text=drawColor,anchor=base,inner sep=0pt, outer sep=0pt, scale=  0.54] at (197.13,189.33) {2.2};

\node[text=drawColor,anchor=base,inner sep=0pt, outer sep=0pt, scale=  0.54] at (197.13, 20.87) {2.6};

\node[text=drawColor,anchor=base,inner sep=0pt, outer sep=0pt, scale=  0.54] at (240.95,124.54) {2.8};

\node[text=drawColor,anchor=base,inner sep=0pt, outer sep=0pt, scale=  0.54] at (240.95, 98.62) {2.9};

\node[text=drawColor,anchor=base,inner sep=0pt, outer sep=0pt, scale=  0.54] at (240.95,202.29) {2.5};

\node[text=drawColor,anchor=base,inner sep=0pt, outer sep=0pt, scale=  0.54] at (240.95,228.21) {1.6};

\node[text=drawColor,anchor=base,inner sep=0pt, outer sep=0pt, scale=  0.54] at (240.95,  7.91) {3.6};

\node[text=drawColor,anchor=base,inner sep=0pt, outer sep=0pt, scale=  0.54] at (240.95,318.92) {1.2};

\node[text=drawColor,anchor=base,inner sep=0pt, outer sep=0pt, scale=  0.54] at (240.95,293.00) {1.4};

\node[text=drawColor,anchor=base,inner sep=0pt, outer sep=0pt, scale=  0.54] at (240.95, 72.71) {3.3};

\node[text=drawColor,anchor=base,inner sep=0pt, outer sep=0pt, scale=  0.54] at (240.95, 33.83) {3.5};

\node[text=drawColor,anchor=base,inner sep=0pt, outer sep=0pt, scale=  0.54] at (240.95,280.04) {1.6};

\node[text=drawColor,anchor=base,inner sep=0pt, outer sep=0pt, scale=  0.54] at (240.95,305.96) {1.2};

\node[text=drawColor,anchor=base,inner sep=0pt, outer sep=0pt, scale=  0.54] at (240.95,137.50) {2.6};

\node[text=drawColor,anchor=base,inner sep=0pt, outer sep=0pt, scale=  0.54] at (240.95, 46.79) {3.5};

\node[text=drawColor,anchor=base,inner sep=0pt, outer sep=0pt, scale=  0.54] at (240.95,254.12) {1.6};

\node[text=drawColor,anchor=base,inner sep=0pt, outer sep=0pt, scale=  0.54] at (240.95, 85.66) {3.2};

\node[text=drawColor,anchor=base,inner sep=0pt, outer sep=0pt, scale=  0.54] at (240.95, 59.75) {3.4};

\node[text=drawColor,anchor=base,inner sep=0pt, outer sep=0pt, scale=  0.54] at (240.95,150.46) {2.7};

\node[text=drawColor,anchor=base,inner sep=0pt, outer sep=0pt, scale=  0.54] at (240.95,241.17) {2.2};

\node[text=drawColor,anchor=base,inner sep=0pt, outer sep=0pt, scale=  0.54] at (240.95,176.37) {3.0};

\node[text=drawColor,anchor=base,inner sep=0pt, outer sep=0pt, scale=  0.54] at (240.95,267.08) {1.6};

\node[text=drawColor,anchor=base,inner sep=0pt, outer sep=0pt, scale=  0.54] at (240.95,215.25) {2.3};

\node[text=drawColor,anchor=base,inner sep=0pt, outer sep=0pt, scale=  0.54] at (240.95,111.58) {2.9};

\node[text=drawColor,anchor=base,inner sep=0pt, outer sep=0pt, scale=  0.54] at (240.95,163.42) {2.6};

\node[text=drawColor,anchor=base,inner sep=0pt, outer sep=0pt, scale=  0.54] at (240.95,189.33) {2.5};

\node[text=drawColor,anchor=base,inner sep=0pt, outer sep=0pt, scale=  0.54] at (240.95, 20.87) {3.3};
\end{scope}
\begin{scope}
\path[clip] (  0.00,  0.00) rectangle (397.48,361.35);
\definecolor{drawColor}{gray}{0.30}

\node[text=drawColor,rotate= 35.00,anchor=base west,inner sep=0pt, outer sep=0pt, scale=  0.70] at ( 68.44,355.40) {Maverick};

\node[text=drawColor,rotate= 35.00,anchor=base west,inner sep=0pt, outer sep=0pt, scale=  0.70] at (112.26,355.40) {Llama 3.3};

\node[text=drawColor,rotate= 35.00,anchor=base west,inner sep=0pt, outer sep=0pt, scale=  0.70] at (156.08,355.40) {Mistral};

\node[text=drawColor,rotate= 35.00,anchor=base west,inner sep=0pt, outer sep=0pt, scale=  0.70] at (199.90,355.40) {Nemotron};

\node[text=drawColor,rotate= 35.00,anchor=base west,inner sep=0pt, outer sep=0pt, scale=  0.70] at (243.71,355.40) {Nova Pro};

\node[text=drawColor,rotate= 35.00,anchor=base west,inner sep=0pt, outer sep=0pt, scale=  0.70] at (287.53,355.40) {Qwen3 235B};

\node[text=drawColor,rotate= 35.00,anchor=base west,inner sep=0pt, outer sep=0pt, scale=  0.70] at (331.35,355.40) {Qwen3 32B};
\end{scope}
\begin{scope}
\path[clip] (  0.00,  0.00) rectangle (397.48,361.35);
\definecolor{drawColor}{gray}{0.30}

\node[text=drawColor,anchor=base east,inner sep=0pt, outer sep=0pt, scale=  0.70] at ( 35.79,  7.36) {Greek};

\node[text=drawColor,anchor=base east,inner sep=0pt, outer sep=0pt, scale=  0.70] at ( 35.79, 20.32) {Ukrainian};

\node[text=drawColor,anchor=base east,inner sep=0pt, outer sep=0pt, scale=  0.70] at ( 35.79, 33.28) {Finnish};

\node[text=drawColor,anchor=base east,inner sep=0pt, outer sep=0pt, scale=  0.70] at ( 35.79, 46.24) {Hungarian};

\node[text=drawColor,anchor=base east,inner sep=0pt, outer sep=0pt, scale=  0.70] at ( 35.79, 59.20) {Latvian};

\node[text=drawColor,anchor=base east,inner sep=0pt, outer sep=0pt, scale=  0.70] at ( 35.79, 72.16) {Estonian};

\node[text=drawColor,anchor=base east,inner sep=0pt, outer sep=0pt, scale=  0.70] at ( 35.79, 85.12) {Lithuanian};

\node[text=drawColor,anchor=base east,inner sep=0pt, outer sep=0pt, scale=  0.70] at ( 35.79, 98.07) {Czech};

\node[text=drawColor,anchor=base east,inner sep=0pt, outer sep=0pt, scale=  0.70] at ( 35.79,111.03) {Slovak};

\node[text=drawColor,anchor=base east,inner sep=0pt, outer sep=0pt, scale=  0.70] at ( 35.79,123.99) {Bulgarian};

\node[text=drawColor,anchor=base east,inner sep=0pt, outer sep=0pt, scale=  0.70] at ( 35.79,136.95) {Croatian};

\node[text=drawColor,anchor=base east,inner sep=0pt, outer sep=0pt, scale=  0.70] at ( 35.79,149.91) {Maltese};

\node[text=drawColor,anchor=base east,inner sep=0pt, outer sep=0pt, scale=  0.70] at ( 35.79,162.87) {Slovenian};

\node[text=drawColor,anchor=base east,inner sep=0pt, outer sep=0pt, scale=  0.70] at ( 35.79,175.82) {Polish};

\node[text=drawColor,anchor=base east,inner sep=0pt, outer sep=0pt, scale=  0.70] at ( 35.79,188.78) {Swedish};

\node[text=drawColor,anchor=base east,inner sep=0pt, outer sep=0pt, scale=  0.70] at ( 35.79,201.74) {Danish};

\node[text=drawColor,anchor=base east,inner sep=0pt, outer sep=0pt, scale=  0.70] at ( 35.79,214.70) {Romanian};

\node[text=drawColor,anchor=base east,inner sep=0pt, outer sep=0pt, scale=  0.70] at ( 35.79,227.66) {German};

\node[text=drawColor,anchor=base east,inner sep=0pt, outer sep=0pt, scale=  0.70] at ( 35.79,240.62) {Dutch};

\node[text=drawColor,anchor=base east,inner sep=0pt, outer sep=0pt, scale=  0.70] at ( 35.79,253.58) {Italian};

\node[text=drawColor,anchor=base east,inner sep=0pt, outer sep=0pt, scale=  0.70] at ( 35.79,266.53) {Portuguese};

\node[text=drawColor,anchor=base east,inner sep=0pt, outer sep=0pt, scale=  0.70] at ( 35.79,279.49) {French};

\node[text=drawColor,anchor=base east,inner sep=0pt, outer sep=0pt, scale=  0.70] at ( 35.79,292.45) {Spanish};

\node[text=drawColor,anchor=base east,inner sep=0pt, outer sep=0pt, scale=  0.70] at ( 35.79,305.41) {Irish};

\node[text=drawColor,anchor=base east,inner sep=0pt, outer sep=0pt, scale=  0.70] at ( 35.79,318.37) {English};
\end{scope}
\begin{scope}
\path[clip] (  0.00,  0.00) rectangle (397.48,361.35);
\definecolor{drawColor}{RGB}{0,0,0}

\node[text=drawColor,anchor=base west,inner sep=0pt, outer sep=0pt, scale=  0.70] at (366.87,221.77) {Fertility};
\end{scope}
\begin{scope}
\path[clip] (  0.00,  0.00) rectangle (397.48,361.35);
\definecolor{fillColor}{RGB}{255,255,178}

\path[fill=fillColor] (366.87,103.28) rectangle (375.41,104.42);
\definecolor{fillColor}{RGB}{255,253,175}

\path[fill=fillColor] (366.87,104.42) rectangle (375.41,105.56);
\definecolor{fillColor}{RGB}{255,251,171}

\path[fill=fillColor] (366.87,105.56) rectangle (375.41,106.69);
\definecolor{fillColor}{RGB}{255,249,168}

\path[fill=fillColor] (366.87,106.69) rectangle (375.41,107.83);
\definecolor{fillColor}{RGB}{255,247,164}

\path[fill=fillColor] (366.87,107.83) rectangle (375.41,108.97);
\definecolor{fillColor}{RGB}{255,245,161}

\path[fill=fillColor] (366.87,108.97) rectangle (375.41,110.11);
\definecolor{fillColor}{RGB}{255,243,157}

\path[fill=fillColor] (366.87,110.11) rectangle (375.41,111.25);
\definecolor{fillColor}{RGB}{255,240,154}

\path[fill=fillColor] (366.87,111.25) rectangle (375.41,112.39);
\definecolor{fillColor}{RGB}{255,238,150}

\path[fill=fillColor] (366.87,112.39) rectangle (375.41,113.52);
\definecolor{fillColor}{RGB}{255,236,147}

\path[fill=fillColor] (366.87,113.52) rectangle (375.41,114.66);
\definecolor{fillColor}{RGB}{255,234,144}

\path[fill=fillColor] (366.87,114.66) rectangle (375.41,115.80);
\definecolor{fillColor}{RGB}{255,232,140}

\path[fill=fillColor] (366.87,115.80) rectangle (375.41,116.94);
\definecolor{fillColor}{RGB}{255,230,137}

\path[fill=fillColor] (366.87,116.94) rectangle (375.41,118.08);
\definecolor{fillColor}{RGB}{255,228,133}

\path[fill=fillColor] (366.87,118.08) rectangle (375.41,119.21);
\definecolor{fillColor}{RGB}{255,226,130}

\path[fill=fillColor] (366.87,119.21) rectangle (375.41,120.35);
\definecolor{fillColor}{RGB}{255,224,126}

\path[fill=fillColor] (366.87,120.35) rectangle (375.41,121.49);
\definecolor{fillColor}{RGB}{255,222,123}

\path[fill=fillColor] (366.87,121.49) rectangle (375.41,122.63);
\definecolor{fillColor}{RGB}{255,220,119}

\path[fill=fillColor] (366.87,122.63) rectangle (375.41,123.77);
\definecolor{fillColor}{RGB}{255,218,116}

\path[fill=fillColor] (366.87,123.77) rectangle (375.41,124.90);
\definecolor{fillColor}{RGB}{255,216,112}

\path[fill=fillColor] (366.87,124.90) rectangle (375.41,126.04);
\definecolor{fillColor}{RGB}{255,214,109}

\path[fill=fillColor] (366.87,126.04) rectangle (375.41,127.18);
\definecolor{fillColor}{RGB}{255,212,105}

\path[fill=fillColor] (366.87,127.18) rectangle (375.41,128.32);
\definecolor{fillColor}{RGB}{255,210,102}

\path[fill=fillColor] (366.87,128.32) rectangle (375.41,129.46);
\definecolor{fillColor}{RGB}{255,208,98}

\path[fill=fillColor] (366.87,129.46) rectangle (375.41,130.59);
\definecolor{fillColor}{RGB}{254,206,95}

\path[fill=fillColor] (366.87,130.59) rectangle (375.41,131.73);
\definecolor{fillColor}{RGB}{254,203,92}

\path[fill=fillColor] (366.87,131.73) rectangle (375.41,132.87);
\definecolor{fillColor}{RGB}{254,201,90}

\path[fill=fillColor] (366.87,132.87) rectangle (375.41,134.01);
\definecolor{fillColor}{RGB}{254,198,89}

\path[fill=fillColor] (366.87,134.01) rectangle (375.41,135.15);
\definecolor{fillColor}{RGB}{254,196,88}

\path[fill=fillColor] (366.87,135.15) rectangle (375.41,136.29);
\definecolor{fillColor}{RGB}{254,194,86}

\path[fill=fillColor] (366.87,136.29) rectangle (375.41,137.42);
\definecolor{fillColor}{RGB}{255,191,85}

\path[fill=fillColor] (366.87,137.42) rectangle (375.41,138.56);
\definecolor{fillColor}{RGB}{255,189,84}

\path[fill=fillColor] (366.87,138.56) rectangle (375.41,139.70);
\definecolor{fillColor}{RGB}{255,186,82}

\path[fill=fillColor] (366.87,139.70) rectangle (375.41,140.84);
\definecolor{fillColor}{RGB}{255,184,81}

\path[fill=fillColor] (366.87,140.84) rectangle (375.41,141.98);
\definecolor{fillColor}{RGB}{255,181,80}

\path[fill=fillColor] (366.87,141.98) rectangle (375.41,143.11);
\definecolor{fillColor}{RGB}{255,179,78}

\path[fill=fillColor] (366.87,143.11) rectangle (375.41,144.25);
\definecolor{fillColor}{RGB}{255,176,77}

\path[fill=fillColor] (366.87,144.25) rectangle (375.41,145.39);
\definecolor{fillColor}{RGB}{255,174,76}

\path[fill=fillColor] (366.87,145.39) rectangle (375.41,146.53);
\definecolor{fillColor}{RGB}{255,171,75}

\path[fill=fillColor] (366.87,146.53) rectangle (375.41,147.67);
\definecolor{fillColor}{RGB}{254,168,73}

\path[fill=fillColor] (366.87,147.67) rectangle (375.41,148.80);
\definecolor{fillColor}{RGB}{254,166,72}

\path[fill=fillColor] (366.87,148.80) rectangle (375.41,149.94);
\definecolor{fillColor}{RGB}{254,163,71}

\path[fill=fillColor] (366.87,149.94) rectangle (375.41,151.08);
\definecolor{fillColor}{RGB}{254,161,69}

\path[fill=fillColor] (366.87,151.08) rectangle (375.41,152.22);
\definecolor{fillColor}{RGB}{254,158,68}

\path[fill=fillColor] (366.87,152.22) rectangle (375.41,153.36);
\definecolor{fillColor}{RGB}{254,156,67}

\path[fill=fillColor] (366.87,153.36) rectangle (375.41,154.50);
\definecolor{fillColor}{RGB}{254,153,66}

\path[fill=fillColor] (366.87,154.50) rectangle (375.41,155.63);
\definecolor{fillColor}{RGB}{254,150,64}

\path[fill=fillColor] (366.87,155.63) rectangle (375.41,156.77);
\definecolor{fillColor}{RGB}{253,148,63}

\path[fill=fillColor] (366.87,156.77) rectangle (375.41,157.91);
\definecolor{fillColor}{RGB}{253,145,62}

\path[fill=fillColor] (366.87,157.91) rectangle (375.41,159.05);
\definecolor{fillColor}{RGB}{253,142,61}

\path[fill=fillColor] (366.87,159.05) rectangle (375.41,160.19);
\definecolor{fillColor}{RGB}{253,140,59}

\path[fill=fillColor] (366.87,160.19) rectangle (375.41,161.32);
\definecolor{fillColor}{RGB}{252,137,58}

\path[fill=fillColor] (366.87,161.32) rectangle (375.41,162.46);
\definecolor{fillColor}{RGB}{252,134,57}

\path[fill=fillColor] (366.87,162.46) rectangle (375.41,163.60);
\definecolor{fillColor}{RGB}{252,131,56}

\path[fill=fillColor] (366.87,163.60) rectangle (375.41,164.74);
\definecolor{fillColor}{RGB}{251,128,55}

\path[fill=fillColor] (366.87,164.74) rectangle (375.41,165.88);
\definecolor{fillColor}{RGB}{251,125,53}

\path[fill=fillColor] (366.87,165.88) rectangle (375.41,167.01);
\definecolor{fillColor}{RGB}{250,122,52}

\path[fill=fillColor] (366.87,167.01) rectangle (375.41,168.15);
\definecolor{fillColor}{RGB}{250,120,51}

\path[fill=fillColor] (366.87,168.15) rectangle (375.41,169.29);
\definecolor{fillColor}{RGB}{249,117,50}

\path[fill=fillColor] (366.87,169.29) rectangle (375.41,170.43);
\definecolor{fillColor}{RGB}{249,114,49}

\path[fill=fillColor] (366.87,170.43) rectangle (375.41,171.57);
\definecolor{fillColor}{RGB}{248,111,48}

\path[fill=fillColor] (366.87,171.57) rectangle (375.41,172.70);
\definecolor{fillColor}{RGB}{248,107,46}

\path[fill=fillColor] (366.87,172.70) rectangle (375.41,173.84);
\definecolor{fillColor}{RGB}{247,104,45}

\path[fill=fillColor] (366.87,173.84) rectangle (375.41,174.98);
\definecolor{fillColor}{RGB}{247,101,44}

\path[fill=fillColor] (366.87,174.98) rectangle (375.41,176.12);
\definecolor{fillColor}{RGB}{246,98,43}

\path[fill=fillColor] (366.87,176.12) rectangle (375.41,177.26);
\definecolor{fillColor}{RGB}{246,95,42}

\path[fill=fillColor] (366.87,177.26) rectangle (375.41,178.40);
\definecolor{fillColor}{RGB}{245,91,41}

\path[fill=fillColor] (366.87,178.40) rectangle (375.41,179.53);
\definecolor{fillColor}{RGB}{245,88,40}

\path[fill=fillColor] (366.87,179.53) rectangle (375.41,180.67);
\definecolor{fillColor}{RGB}{244,84,39}

\path[fill=fillColor] (366.87,180.67) rectangle (375.41,181.81);
\definecolor{fillColor}{RGB}{243,81,38}

\path[fill=fillColor] (366.87,181.81) rectangle (375.41,182.95);
\definecolor{fillColor}{RGB}{243,77,36}

\path[fill=fillColor] (366.87,182.95) rectangle (375.41,184.09);
\definecolor{fillColor}{RGB}{242,73,35}

\path[fill=fillColor] (366.87,184.09) rectangle (375.41,185.22);
\definecolor{fillColor}{RGB}{241,69,34}

\path[fill=fillColor] (366.87,185.22) rectangle (375.41,186.36);
\definecolor{fillColor}{RGB}{241,65,33}

\path[fill=fillColor] (366.87,186.36) rectangle (375.41,187.50);
\definecolor{fillColor}{RGB}{240,60,32}

\path[fill=fillColor] (366.87,187.50) rectangle (375.41,188.64);
\definecolor{fillColor}{RGB}{238,58,32}

\path[fill=fillColor] (366.87,188.64) rectangle (375.41,189.78);
\definecolor{fillColor}{RGB}{236,56,33}

\path[fill=fillColor] (366.87,189.78) rectangle (375.41,190.91);
\definecolor{fillColor}{RGB}{234,54,33}

\path[fill=fillColor] (366.87,190.91) rectangle (375.41,192.05);
\definecolor{fillColor}{RGB}{232,53,33}

\path[fill=fillColor] (366.87,192.05) rectangle (375.41,193.19);
\definecolor{fillColor}{RGB}{230,51,34}

\path[fill=fillColor] (366.87,193.19) rectangle (375.41,194.33);
\definecolor{fillColor}{RGB}{228,49,34}

\path[fill=fillColor] (366.87,194.33) rectangle (375.41,195.47);
\definecolor{fillColor}{RGB}{226,47,34}

\path[fill=fillColor] (366.87,195.47) rectangle (375.41,196.61);
\definecolor{fillColor}{RGB}{224,45,35}

\path[fill=fillColor] (366.87,196.61) rectangle (375.41,197.74);
\definecolor{fillColor}{RGB}{222,43,35}

\path[fill=fillColor] (366.87,197.74) rectangle (375.41,198.88);
\definecolor{fillColor}{RGB}{220,42,35}

\path[fill=fillColor] (366.87,198.88) rectangle (375.41,200.02);
\definecolor{fillColor}{RGB}{218,40,36}

\path[fill=fillColor] (366.87,200.02) rectangle (375.41,201.16);
\definecolor{fillColor}{RGB}{216,38,36}

\path[fill=fillColor] (366.87,201.16) rectangle (375.41,202.30);
\definecolor{fillColor}{RGB}{214,36,36}

\path[fill=fillColor] (366.87,202.30) rectangle (375.41,203.43);
\definecolor{fillColor}{RGB}{211,33,36}

\path[fill=fillColor] (366.87,203.43) rectangle (375.41,204.57);
\definecolor{fillColor}{RGB}{209,31,36}

\path[fill=fillColor] (366.87,204.57) rectangle (375.41,205.71);
\definecolor{fillColor}{RGB}{207,29,37}

\path[fill=fillColor] (366.87,205.71) rectangle (375.41,206.85);
\definecolor{fillColor}{RGB}{205,27,37}

\path[fill=fillColor] (366.87,206.85) rectangle (375.41,207.99);
\definecolor{fillColor}{RGB}{203,24,37}

\path[fill=fillColor] (366.87,207.99) rectangle (375.41,209.12);
\definecolor{fillColor}{RGB}{201,22,37}

\path[fill=fillColor] (366.87,209.12) rectangle (375.41,210.26);
\definecolor{fillColor}{RGB}{199,19,37}

\path[fill=fillColor] (366.87,210.26) rectangle (375.41,211.40);
\definecolor{fillColor}{RGB}{197,16,37}

\path[fill=fillColor] (366.87,211.40) rectangle (375.41,212.54);
\definecolor{fillColor}{RGB}{195,12,38}

\path[fill=fillColor] (366.87,212.54) rectangle (375.41,213.68);
\definecolor{fillColor}{RGB}{193,8,38}

\path[fill=fillColor] (366.87,213.68) rectangle (375.41,214.82);
\definecolor{fillColor}{RGB}{191,4,38}

\path[fill=fillColor] (366.87,214.82) rectangle (375.41,215.95);
\definecolor{fillColor}{RGB}{189,0,38}

\path[fill=fillColor] (366.87,215.95) rectangle (375.41,217.09);
\end{scope}
\begin{scope}
\path[clip] (  0.00,  0.00) rectangle (397.48,361.35);
\definecolor{drawColor}{RGB}{255,255,255}

\path[draw=drawColor,line width= 0.2pt,line join=round] (372.52,103.85) --
	(375.41,103.85);

\path[draw=drawColor,line width= 0.2pt,line join=round] (372.52,127.82) --
	(375.41,127.82);

\path[draw=drawColor,line width= 0.2pt,line join=round] (372.52,151.80) --
	(375.41,151.80);

\path[draw=drawColor,line width= 0.2pt,line join=round] (372.52,175.77) --
	(375.41,175.77);

\path[draw=drawColor,line width= 0.2pt,line join=round] (372.52,199.74) --
	(375.41,199.74);

\path[draw=drawColor,line width= 0.2pt,line join=round] (369.77,103.85) --
	(366.87,103.85);

\path[draw=drawColor,line width= 0.2pt,line join=round] (369.77,127.82) --
	(366.87,127.82);

\path[draw=drawColor,line width= 0.2pt,line join=round] (369.77,151.80) --
	(366.87,151.80);

\path[draw=drawColor,line width= 0.2pt,line join=round] (369.77,175.77) --
	(366.87,175.77);

\path[draw=drawColor,line width= 0.2pt,line join=round] (369.77,199.74) --
	(366.87,199.74);
\end{scope}
\begin{scope}
\path[clip] (  0.00,  0.00) rectangle (397.48,361.35);
\definecolor{drawColor}{RGB}{0,0,0}

\node[text=drawColor,anchor=base west,inner sep=0pt, outer sep=0pt, scale=  0.60] at (379.41,101.78) {1};

\node[text=drawColor,anchor=base west,inner sep=0pt, outer sep=0pt, scale=  0.60] at (379.41,125.76) {2};

\node[text=drawColor,anchor=base west,inner sep=0pt, outer sep=0pt, scale=  0.60] at (379.41,149.73) {3};

\node[text=drawColor,anchor=base west,inner sep=0pt, outer sep=0pt, scale=  0.60] at (379.41,173.70) {4};

\node[text=drawColor,anchor=base west,inner sep=0pt, outer sep=0pt, scale=  0.60] at (379.41,197.67) {5};
\end{scope}
\end{tikzpicture}

\caption{Tokenizer fertility across 25~languages and 7~models on identical EU legislative text. Languages sorted by mean fertility. Ukrainian (3.75 mean) is more expensive than all other Slavic languages. Greek (5.66 mean) pays the steepest penalty due to script divergence.}
\label{fig:eu_heatmap}
\end{figure*}

Figure~\ref{fig:eu_heatmap} shows the full heatmap across all seven models and 25~languages. The hierarchy is consistent across models: English and Irish are always cheapest; Greek and Ukrainian always most expensive after Greek. Model rankings are also stable: both Qwen models are consistently the least efficient across all languages, while Maverick is the most efficient.

\subsection{Experiment 3: Cross-Lingual Few-Shot Classification}
\label{sec:exp3}

To test whether the few-shot effect is language-specific or model-intrinsic, we evaluated all seven models on SIB-200 topic classification across four Slavic languages using parallel examples (same \texttt{index\_id}). We also include the legal text few-shot delta from \citet{ovcharov2026tokenizer} for cross-domain comparison.

\begin{table*}[t]
\centering
\caption{Few-shot delta ($\Delta_{\text{FS}}$ = FS $-$ ZS accuracy, in pp) on SIB-200 topic classification across four Slavic languages, plus Ukrainian legal text from \citet{ovcharov2026tokenizer}. Models sorted by SIB-200 average. Bold = consistent direction across all 4 news languages.}
\label{tab:slavic}
\begin{tabular}{lR{1.0cm}R{1.0cm}R{1.0cm}R{1.0cm}R{1.0cm}R{1.4cm}}
\toprule
& \multicolumn{5}{c}{\textbf{SIB-200 (news)}} & \textbf{Legal} \\
\cmidrule(lr){2-6}
\textbf{Model} & \textbf{UK} & \textbf{PL} & \textbf{RU} & \textbf{CZ} & \textbf{Avg} & \textbf{UK} \\
\midrule
\textbf{Nemotron Super 3}  & $+$8.3  & $+$15.7 & $+$8.8  & $+$9.3  & \textbf{$+$10.5} & $-$12.8 \\
Qwen3 32B         & $+$3.9  & $+$5.9  & $+$7.4  & $+$7.8  & \textbf{$+$6.2} & $-$6.6 \\
Nova Pro          & $+$4.4  & $+$5.1  & $+$3.0  & $+$1.5  & \textbf{$+$3.5} & $-$2.6 \\
Qwen3 235B        & $+$4.9  & $+$3.4  & $+$1.5  & $+$2.9  & \textbf{$+$3.2} & $-$26.0 \\
Llama 3.3 70B     & $+$3.4  & $+$3.0  & $-$0.5  & $+$5.0  & $+$2.7 & $+$0.4 \\
Mistral Large 3   & $-$4.4  & $-$1.0  & $+$2.9  & $+$1.0  & $-$0.4 & $-$3.3 \\
\textbf{Llama 4 Maverick}  & $-$7.8  & $-$6.3  & $-$15.7 & $-$3.8  & \textbf{$-$8.4} & $-$6.2 \\
\bottomrule
\end{tabular}
\end{table*}

With 204 test examples and 7~classes, the 95\% Clopper--Pearson confidence interval on a single accuracy is $\pm$5--7~pp. To assess whether few-shot deltas exceed this noise floor, we compute exact McNemar's tests on paired zero-shot vs.\ few-shot predictions for each model$\times$language cell. Of the 28~cells (7~models $\times$ 4~languages), 12~are significant at $p < 0.05$ after Holm--Bonferroni correction: all four Maverick cells (degradation) and all four Nemotron cells (improvement), plus Qwen3~32B on Russian and Czech, and Llama~3.3 on Czech. The remaining deltas, while directionally consistent, do not individually reach significance at this sample size.

The results answer the central question of this paper. The few-shot effect pattern is \textbf{identical across all four Slavic languages}:

\begin{itemize}[leftmargin=*]
    \item Four models (Nemotron, Qwen3~32B, Qwen3~235B, Nova~Pro) \textbf{improve on all four languages}---average deltas range from $+$3.2 to $+$10.5~pp. Nemotron's improvement is significant on all four ($p < 0.01$).
    \item One model (Llama~4 Maverick) \textbf{degrades on all four languages}---average $-$8.4~pp, with Russian showing the most severe degradation ($-$15.7~pp). All four Maverick cells are significant ($p < 0.001$).
    \item Mistral Large~3 and Llama~3.3 show mixed effects, with language-dependent direction; individual deltas are not significant.
\end{itemize}

By language, 5--6 of 7 models improve with few-shot examples across all four languages (Czech: 6/7, others: 5/7). The consistency across languages with different scripts (Cyrillic for Ukrainian and Russian vs.\ Latin for Polish and Czech), different tokenizer fertility levels, and different pre-training data volumes rules out both the ``Ukrainian-specific'' and ``morphological complexity'' hypotheses.

The few-shot effect is \textbf{model-intrinsic}: Maverick's attention mechanism consistently anchors on surface patterns in demonstrations regardless of language, while Nemotron's hybrid architecture consistently leverages demonstrations for task specification. This is a property of model design, not of the input language.

\paragraph{Error analysis: Maverick degradation.} To understand \emph{how} Maverick degrades, we examined its Ukrainian SIB-200 predictions. In zero-shot mode, errors are distributed across classes roughly proportionally to class frequency. In few-shot mode, Maverick shifts toward two dominant classes---\emph{science\_and\_technology} and \emph{politics}---at the expense of minority classes: \emph{geography} recall drops from 76.5\% to 41.2\%, and \emph{entertainment} from 82.1\% to 64.3\%. This pattern is consistent with \emph{demonstration anchoring}: the model over-weights surface patterns from the provided examples (which include one geography and one entertainment example) rather than using them for task specification. The same class-collapse pattern appears on Russian ($-$15.7~pp), where Maverick's geography recall drops from 70.6\% to 23.5\%. \citet{ovcharov2026fewshot} provide a mechanistic explanation: in models vulnerable to few-shot degradation, hidden-state representations shift toward the demonstration content rather than the query, a phenomenon quantified by their \emph{content delta} metric.

\begin{figure*}[t]
\centering
% Created by tikzDevice version 0.12.6 on 2026-05-15 12:03:01
% !TEX encoding = UTF-8 Unicode
\begin{tikzpicture}[x=1pt,y=1pt]
\definecolor{fillColor}{RGB}{255,255,255}
\path[use as bounding box,fill=fillColor,fill opacity=0.00] (0,0) rectangle (397.48,216.81);
\begin{scope}
\path[clip] (  0.00,  0.00) rectangle (397.48,216.81);
\definecolor{fillColor}{RGB}{255,255,255}

\path[fill=fillColor] (  0.00,  0.00) rectangle (397.48,216.81);
\end{scope}
\begin{scope}
\path[clip] ( 44.19, 19.96) rectangle (128.01,201.33);
\definecolor{drawColor}{gray}{0.92}

\path[draw=drawColor,line width= 0.4pt,line join=round] ( 61.83, 19.96) --
	( 61.83,201.33);

\path[draw=drawColor,line width= 0.4pt,line join=round] ( 86.10, 19.96) --
	( 86.10,201.33);

\path[draw=drawColor,line width= 0.4pt,line join=round] (110.37, 19.96) --
	(110.37,201.33);
\definecolor{fillColor}{RGB}{192,57,43}

\path[fill=fillColor] ( 67.17, 26.26) rectangle ( 86.10, 43.89);
\definecolor{fillColor}{RGB}{39,174,96}

\path[fill=fillColor] ( 86.10, 76.64) rectangle ( 94.35, 94.28);
\definecolor{fillColor}{RGB}{192,57,43}

\path[fill=fillColor] ( 75.42, 51.45) rectangle ( 86.10, 69.08);
\definecolor{fillColor}{RGB}{39,174,96}

\path[fill=fillColor] ( 86.10,177.40) rectangle (106.24,195.04);

\path[fill=fillColor] ( 86.10,127.02) rectangle ( 96.78,144.66);

\path[fill=fillColor] ( 86.10,101.83) rectangle ( 97.99,119.47);

\path[fill=fillColor] ( 86.10,152.21) rectangle ( 95.56,169.85);
\definecolor{drawColor}{gray}{0.30}

\path[draw=drawColor,line width= 0.3pt,line join=round] ( 86.10, 19.96) -- ( 86.10,201.33);
\end{scope}
\begin{scope}
\path[clip] (132.01, 19.96) rectangle (215.84,201.33);
\definecolor{drawColor}{gray}{0.92}

\path[draw=drawColor,line width= 0.4pt,line join=round] (149.66, 19.96) --
	(149.66,201.33);

\path[draw=drawColor,line width= 0.4pt,line join=round] (173.92, 19.96) --
	(173.92,201.33);

\path[draw=drawColor,line width= 0.4pt,line join=round] (198.19, 19.96) --
	(198.19,201.33);
\definecolor{fillColor}{RGB}{192,57,43}

\path[fill=fillColor] (158.63, 26.26) rectangle (173.92, 43.89);
\definecolor{fillColor}{RGB}{39,174,96}

\path[fill=fillColor] (173.92, 76.64) rectangle (181.20, 94.28);
\definecolor{fillColor}{RGB}{192,57,43}

\path[fill=fillColor] (171.50, 51.45) rectangle (173.92, 69.08);
\definecolor{fillColor}{RGB}{39,174,96}

\path[fill=fillColor] (173.92,177.40) rectangle (212.03,195.04);

\path[fill=fillColor] (173.92,127.02) rectangle (186.30,144.66);

\path[fill=fillColor] (173.92,101.83) rectangle (182.18,119.47);

\path[fill=fillColor] (173.92,152.21) rectangle (188.24,169.85);
\definecolor{drawColor}{gray}{0.30}

\path[draw=drawColor,line width= 0.3pt,line join=round] (173.92, 19.96) -- (173.92,201.33);
\end{scope}
\begin{scope}
\path[clip] (219.84, 19.96) rectangle (303.66,201.33);
\definecolor{drawColor}{gray}{0.92}

\path[draw=drawColor,line width= 0.4pt,line join=round] (237.48, 19.96) --
	(237.48,201.33);

\path[draw=drawColor,line width= 0.4pt,line join=round] (261.75, 19.96) --
	(261.75,201.33);

\path[draw=drawColor,line width= 0.4pt,line join=round] (286.02, 19.96) --
	(286.02,201.33);
\definecolor{fillColor}{RGB}{192,57,43}

\path[fill=fillColor] (223.65, 26.26) rectangle (261.75, 43.89);

\path[fill=fillColor] (260.54, 76.64) rectangle (261.75, 94.28);
\definecolor{fillColor}{RGB}{39,174,96}

\path[fill=fillColor] (261.75, 51.45) rectangle (268.79, 69.08);

\path[fill=fillColor] (261.75,177.40) rectangle (283.11,195.04);

\path[fill=fillColor] (261.75,127.02) rectangle (269.03,144.66);

\path[fill=fillColor] (261.75,101.83) rectangle (265.39,119.47);

\path[fill=fillColor] (261.75,152.21) rectangle (279.71,169.85);
\definecolor{drawColor}{gray}{0.30}

\path[draw=drawColor,line width= 0.3pt,line join=round] (261.75, 19.96) -- (261.75,201.33);
\end{scope}
\begin{scope}
\path[clip] (307.66, 19.96) rectangle (391.48,201.33);
\definecolor{drawColor}{gray}{0.92}

\path[draw=drawColor,line width= 0.4pt,line join=round] (325.30, 19.96) --
	(325.30,201.33);

\path[draw=drawColor,line width= 0.4pt,line join=round] (349.57, 19.96) --
	(349.57,201.33);

\path[draw=drawColor,line width= 0.4pt,line join=round] (373.84, 19.96) --
	(373.84,201.33);
\definecolor{fillColor}{RGB}{192,57,43}

\path[fill=fillColor] (340.35, 26.26) rectangle (349.57, 43.89);
\definecolor{fillColor}{RGB}{39,174,96}

\path[fill=fillColor] (349.57, 76.64) rectangle (361.71, 94.28);

\path[fill=fillColor] (349.57, 51.45) rectangle (352.00, 69.08);

\path[fill=fillColor] (349.57,177.40) rectangle (372.14,195.04);

\path[fill=fillColor] (349.57,127.02) rectangle (353.21,144.66);

\path[fill=fillColor] (349.57,101.83) rectangle (356.61,119.47);

\path[fill=fillColor] (349.57,152.21) rectangle (368.50,169.85);
\definecolor{drawColor}{gray}{0.30}

\path[draw=drawColor,line width= 0.3pt,line join=round] (349.57, 19.96) -- (349.57,201.33);
\end{scope}
\begin{scope}
\path[clip] ( 44.19,201.33) rectangle (128.01,214.81);
\definecolor{drawColor}{gray}{0.10}

\node[text=drawColor,anchor=base,inner sep=0pt, outer sep=0pt, scale=  0.80] at ( 86.10,205.31) {\bfseries Ukrainian};
\end{scope}
\begin{scope}
\path[clip] (132.01,201.33) rectangle (215.84,214.81);
\definecolor{drawColor}{gray}{0.10}

\node[text=drawColor,anchor=base,inner sep=0pt, outer sep=0pt, scale=  0.80] at (173.92,205.31) {\bfseries Polish};
\end{scope}
\begin{scope}
\path[clip] (219.84,201.33) rectangle (303.66,214.81);
\definecolor{drawColor}{gray}{0.10}

\node[text=drawColor,anchor=base,inner sep=0pt, outer sep=0pt, scale=  0.80] at (261.75,205.31) {\bfseries Russian};
\end{scope}
\begin{scope}
\path[clip] (307.66,201.33) rectangle (391.48,214.81);
\definecolor{drawColor}{gray}{0.10}

\node[text=drawColor,anchor=base,inner sep=0pt, outer sep=0pt, scale=  0.80] at (349.57,205.31) {\bfseries Czech};
\end{scope}
\begin{scope}
\path[clip] (  0.00,  0.00) rectangle (397.48,216.81);
\definecolor{drawColor}{gray}{0.30}

\node[text=drawColor,anchor=base,inner sep=0pt, outer sep=0pt, scale=  0.60] at ( 61.83, 12.23) {-10};

\node[text=drawColor,anchor=base,inner sep=0pt, outer sep=0pt, scale=  0.60] at ( 86.10, 12.23) {0};

\node[text=drawColor,anchor=base,inner sep=0pt, outer sep=0pt, scale=  0.60] at (110.37, 12.23) {10};
\end{scope}
\begin{scope}
\path[clip] (  0.00,  0.00) rectangle (397.48,216.81);
\definecolor{drawColor}{gray}{0.30}

\node[text=drawColor,anchor=base,inner sep=0pt, outer sep=0pt, scale=  0.60] at (149.66, 12.23) {-10};

\node[text=drawColor,anchor=base,inner sep=0pt, outer sep=0pt, scale=  0.60] at (173.92, 12.23) {0};

\node[text=drawColor,anchor=base,inner sep=0pt, outer sep=0pt, scale=  0.60] at (198.19, 12.23) {10};
\end{scope}
\begin{scope}
\path[clip] (  0.00,  0.00) rectangle (397.48,216.81);
\definecolor{drawColor}{gray}{0.30}

\node[text=drawColor,anchor=base,inner sep=0pt, outer sep=0pt, scale=  0.60] at (237.48, 12.23) {-10};

\node[text=drawColor,anchor=base,inner sep=0pt, outer sep=0pt, scale=  0.60] at (261.75, 12.23) {0};

\node[text=drawColor,anchor=base,inner sep=0pt, outer sep=0pt, scale=  0.60] at (286.02, 12.23) {10};
\end{scope}
\begin{scope}
\path[clip] (  0.00,  0.00) rectangle (397.48,216.81);
\definecolor{drawColor}{gray}{0.30}

\node[text=drawColor,anchor=base,inner sep=0pt, outer sep=0pt, scale=  0.60] at (325.30, 12.23) {-10};

\node[text=drawColor,anchor=base,inner sep=0pt, outer sep=0pt, scale=  0.60] at (349.57, 12.23) {0};

\node[text=drawColor,anchor=base,inner sep=0pt, outer sep=0pt, scale=  0.60] at (373.84, 12.23) {10};
\end{scope}
\begin{scope}
\path[clip] (  0.00,  0.00) rectangle (397.48,216.81);
\definecolor{drawColor}{gray}{0.30}

\node[text=drawColor,anchor=base east,inner sep=0pt, outer sep=0pt, scale=  0.70] at ( 40.59, 32.67) {Maverick};

\node[text=drawColor,anchor=base east,inner sep=0pt, outer sep=0pt, scale=  0.70] at ( 40.59, 57.86) {Mistral};

\node[text=drawColor,anchor=base east,inner sep=0pt, outer sep=0pt, scale=  0.70] at ( 40.59, 83.05) {Llama 3.3};

\node[text=drawColor,anchor=base east,inner sep=0pt, outer sep=0pt, scale=  0.70] at ( 40.59,108.24) {Qwen3 235B};

\node[text=drawColor,anchor=base east,inner sep=0pt, outer sep=0pt, scale=  0.70] at ( 40.59,133.43) {Nova Pro};

\node[text=drawColor,anchor=base east,inner sep=0pt, outer sep=0pt, scale=  0.70] at ( 40.59,158.62) {Qwen3 32B};

\node[text=drawColor,anchor=base east,inner sep=0pt, outer sep=0pt, scale=  0.70] at ( 40.59,183.81) {Nemotron};
\end{scope}
\begin{scope}
\path[clip] (  0.00,  0.00) rectangle (397.48,216.81);
\definecolor{drawColor}{RGB}{0,0,0}

\node[text=drawColor,anchor=base,inner sep=0pt, outer sep=0pt, scale=  0.80] at (217.84,  3.56) {$\Delta_{\mathrm{FS}}$ (pp)};
\end{scope}
\end{tikzpicture}

\caption{Few-shot delta ($\Delta_{\text{FS}}$, pp) across four Slavic languages on SIB-200. Red bars = degradation, green = improvement. Maverick degrades on all four languages; Nemotron improves on all four. The pattern is model-intrinsic, not language-dependent.}
\label{fig:fewshot_delta}
\end{figure*}

\subsection{Experiment 4: Linguistic Competence}
\label{sec:exp4}

To test whether tokenizer fertility predicts broader linguistic capability, we evaluated all seven models on the ULP benchmark---347 expert-curated multiple-choice questions testing Ukrainian grammar and orthography.

\begin{table}[t]
\centering
\caption{ULP Ukrainian Language Proficiency results. Fertility from \citet{ovcharov2026tokenizer}. $\Delta_{\text{FS}}$ = few-shot minus zero-shot accuracy.}
\label{tab:ulp}
\begin{tabular}{lR{1.2cm}R{1.2cm}R{1.2cm}R{1.4cm}}
\toprule
\textbf{Model} & \textbf{Fert.} & \textbf{ZS} & \textbf{FS} & \textbf{$\Delta_{\text{FS}}$} \\
\midrule
Llama 4 Maverick   & 2.43 & 57.3\% & 53.2\% & $-$4.2 \\
Llama 3.3 70B      & 2.65 & 33.4\% & 36.9\% & $+$3.5 \\
Nova Pro           & 2.85 & 42.9\% & 41.6\% & $-$1.4 \\
Mistral Large 3    & 3.06 & 51.0\% & 49.4\% & $-$1.6 \\
Nemotron Super 3   & 3.08 & 27.7\% & 36.3\% & $+$8.7 \\
Qwen3 235B         & 3.89 & 42.1\% & 44.2\% & $+$2.1 \\
Qwen3 32B          & 3.90 & 34.6\% & 31.7\% & $-$2.9 \\
\bottomrule
\end{tabular}
\end{table}

Llama~4 Maverick---the model with the most efficient tokenizer (fertility 2.43)---achieves the highest ULP accuracy (57.3\%), 6~percentage points above the next-best model (Mistral Large~3, 51.0\%). However, the overall correlation between fertility and zero-shot accuracy is weak: Spearman $\rho = -0.43$ ($p = 0.34$), Pearson $r = -0.35$ ($p = 0.44$). With only seven data points, statistical power is limited, but the trend direction is consistent with the hypothesis that better tokenization facilitates linguistic competence.

Nemotron Super~3 is a notable outlier: despite moderate fertility (3.08), it achieves only 27.7\% zero-shot accuracy---the lowest among all models, and barely above the 20--25\% random baseline for 4--5 choice questions. This model's hybrid Mamba-Transformer architecture may be optimized for long-document reasoning rather than fine-grained grammatical knowledge.

The few-shot effect on ULP is mixed: four models degrade, three improve. The largest improvement is Nemotron Super~3 ($+$8.7~pp), partially compensating for its low zero-shot baseline. The largest degradation is Maverick ($-$4.2~pp), consistent with its pattern of few-shot sensitivity observed on SIB-200.

% ============================================================
% 5. DISCUSSION
% ============================================================
\section{Discussion}

\subsection{Fertility Is a Tokenizer Property}

Experiment~1 demonstrates that tokenizer fertility rankings are invariant across domains. The 1.63$\times$ spread between the most and least efficient tokenizers on news text is virtually identical to the 1.61$\times$ spread on legal text, and model rankings are preserved: both Qwen models cluster at the inefficient end (${\sim}$3.58 tokens/word), both Llama models at the efficient end (${\sim}$2.2--2.3), with Mistral and Nemotron in between.

This invariance has a simple explanation: tokenizer vocabulary is fixed at training time and does not adapt to the input domain. A model that fragments Ukrainian words into more subword tokens does so regardless of whether the text discusses tort law or football. The absolute fertility level is lower on news text (6 of 7 models show a negative $\Delta$), reflecting the higher density of domain-specific terminology in legal text, but the \emph{relative ranking} is determined by vocabulary design.

Experiment~2 extends this finding across 24~EU languages: on identical legal content, fertility varies by 4.6$\times$ between the most efficient (English, 1.24) and least efficient (Greek, 5.66) languages. Crucially, within the Slavic family, Ukrainian (3.75) is 20--40\% more expensive than Polish (2.64) or Czech (3.13), suggesting that Ukrainian's penalty combines morphological complexity with underrepresentation in pre-training data. This has direct cost implications: processing the same EU regulation in Ukrainian costs 3$\times$ more than in English, purely due to tokenizer design.

\subsection{Few-Shot Effects Are Model-Intrinsic, Confirmed Cross-Lingually}

Experiment~3 confirms on API models what \citet{ovcharov2026fewshot} demonstrated mechanistically on open-weight models: the few-shot effect is task-dependent, with five of seven models improving on SIB-200 news while five of seven degraded on legal text. Two models -- Maverick ($-$7.8~pp) and Mistral ($-$4.4~pp) -- degrade on both tasks.

The cross-lingual dimension (Experiment~3) adds a finding that single-language studies cannot provide: the effect is not just task-dependent but \textbf{model-intrinsic across languages}. Maverick degrades on all four Slavic languages ($-$8.4~pp average); Nemotron improves on all four ($+$10.5~pp). This cross-lingual consistency, combined with the mechanistic explanation from \citet{ovcharov2026fewshot} (content delta predicting benefit), forms a complete picture: the few-shot effect is determined by how a model processes demonstration content, not by the input language.

\subsection{The Ukrainian Penalty: Morphology Plus Underrepresentation}

Experiment~2 reveals that Ukrainian's tokenizer tax is not explained by morphological complexity alone. Within the Slavic family, Polish (2.64) and Czech (3.13) -- languages with comparable inflectional systems -- are 20--40\% cheaper to tokenize than Ukrainian (3.75). This gap likely reflects differences in digital presence and pre-training data volume.

Publicly available corpus statistics confirm this hypothesis. Table~\ref{tab:pretraining} shows the volume of Ukrainian vs.\ other Slavic languages in three major pre-training corpora. Despite comparable speaker populations (${\sim}$38--40M for Ukrainian and Polish), Ukrainian consistently has 2--6$\times$ less training data: 3.1$\times$ fewer tokens than Polish in CulturaX \citep{nguyen2024culturax}, 2.4$\times$ less text in mC4, and 5.7$\times$ fewer words in OSCAR~23.01 \citep{abadji2022towards}. Czech, with only 10.5M speakers, has 1.2--3.0$\times$ more data than Ukrainian in most corpora. This data deficit directly explains the tokenizer penalty: subword vocabularies are optimized on training data, and languages with less data receive fewer dedicated merge operations, resulting in more fragmented tokenization.

\begin{table}[t]
\centering
\caption{Ukrainian vs.\ other Slavic languages in major pre-training corpora. Ratio = language volume / Ukrainian volume.}
\label{tab:pretraining}
\begin{tabular}{lR{1.4cm}R{1.0cm}R{1.4cm}R{1.0cm}}
\toprule
\textbf{Corpus} & \textbf{UK} & \textbf{PL} & \textbf{CZ} & \textbf{RU} \\
\midrule
mC4 (GB)       & 196   & 473 (2.4$\times$)  & 235 (1.2$\times$)  & 3{,}615 (18$\times$) \\
OSCAR 23.01 (B words) & 3.2  & 18.1 (5.7$\times$) & 9.7 (3.0$\times$)  & 78.0 (24$\times$) \\
CulturaX (B tokens)   & 38.2 & 117.3 (3.1$\times$) & 56.9 (1.5$\times$) & 737.2 (19$\times$) \\
\bottomrule
\end{tabular}
\end{table}

The practical consequence is concrete: processing the same EU regulation in Ukrainian costs 3$\times$ more than in English, purely due to tokenizer design. For production systems serving Ukrainian users, this tax compounds across every API call.

\subsection{Fertility and Linguistic Competence}

Experiment~4 is exploratory: with $n = 7$ models, it cannot establish a robust fertility--competence relationship. The correlation on ULP is suggestive but not statistically significant ($\rho = -0.43$, $p = 0.34$). Maverick (lowest fertility, highest ULP) and both Qwen models (highest fertility, below-average ULP) fit the expected pattern, but Nemotron Super~3 (moderate fertility, worst ULP) and Mistral Large~3 (moderate fertility, second-best ULP) break it.

This suggests that fertility is at best a \emph{weak} predictor of linguistic competence. A well-optimized tokenizer that preserves morphological boundaries may facilitate grammar tasks, but architecture and training data composition matter at least as much. Nemotron's Mamba-Transformer hybrid, optimized for long-range dependencies, appears to trade fine-grained morphological sensitivity for document-level reasoning---excelling at case outcome classification \citep[96.0\%]{ovcharov2026tokenizer} while failing on grammatical minutiae (27.7\% ULP).

\subsection{Practical Recommendations}

\begin{enumerate}[leftmargin=*]
    \item \textbf{Default to zero-shot} for morphologically rich languages. Validate few-shot per model and task before deploying. Consider chain-of-thought \citep{wei2022chain} as an alternative prompting strategy that does not require demonstration examples.
    \item \textbf{Start with tokenizer analysis.} Fertility rankings are domain-invariant; measure once, apply everywhere.
    \item \textbf{Ignore parameter counts} for non-English model selection. Use language-specific benchmarks.
    \item \textbf{Budget for the tokenizer tax.} A model with 1.6$\times$ higher fertility is 1.6$\times$ more expensive per document.
\end{enumerate}

% ============================================================
% 6. LIMITATIONS
% ============================================================
\section{Limitations}

\paragraph{SIB-200 test set size.} The test split contains only 204~examples with 7~classes, resulting in some classes having as few as 17~instances (geography). Confidence intervals on minority classes are wide.

\paragraph{Two domains only.} Domain invariance of fertility is tested on two domains (legal, news). Adding a third domain (e.g., biomedical, parliamentary) would strengthen the invariance claim. Similarly, the ``task-dependent'' few-shot conclusion rests on two tasks; a third classification task would increase confidence.

\paragraph{Slavic language coverage.} We test four Slavic languages. Including non-Slavic morphologically rich languages (e.g., Finnish, Hungarian, Turkish) would strengthen the morphological complexity hypothesis.

\paragraph{TMX segment quality.} EU Act translations may contain formulaic boilerplate that inflates cross-lingual similarity and deflates fertility differences.

\paragraph{ULP dataset size.} With 347~questions and only 7~data points in the fertility--competence correlation, statistical power is limited.

\paragraph{Same models as \citet{ovcharov2026tokenizer}.} We use the same seven models for comparability, but the landscape evolves rapidly. Results may not generalize to models released after May~2026.

% ============================================================
% 7. CONCLUSION
% ============================================================
\section{Conclusion}

We presented the first controlled cross-lingual tokenizer fertility map for 24~European languages and tested its downstream consequences. Our key findings:

\begin{enumerate}[leftmargin=*]
    \item \textbf{The tokenizer tax spans 4.6$\times$} across European languages on identical legal text, from English (1.24) to Greek (5.66). Ukrainian (3.75) pays 20--40\% more than cognate Slavic languages, combining morphological complexity with pre-training underrepresentation.

    \item \textbf{Fertility is domain-invariant.} The 1.6$\times$ spread between models is identical on legal (1.61$\times$) and news (1.63$\times$) text. A single measurement predicts cost across all domains.

    \item \textbf{Few-shot effects are model-intrinsic across languages.} Cross-lingual experiments on four Slavic languages show identical patterns: Maverick degrades everywhere ($-$8.4~pp avg), Nemotron improves everywhere ($+$10.5~pp avg), confirming the mechanistic findings of \citet{ovcharov2026fewshot}.

    \item \textbf{Fertility weakly predicts grammatical competence} ($\rho = -0.43$, $n = 7$), but architecture matters more: Nemotron excels at classification while scoring worst on grammar.
\end{enumerate}

For practitioners: \textbf{measure fertility once} -- it predicts cost across all domains. \textbf{Validate few-shot per model and task}, not per language. \textbf{Budget for the Ukrainian tax}: 3$\times$ English cost per API call. Total experiment cost: \$4.05.

% ============================================================
% ACKNOWLEDGMENTS
% ============================================================
\section*{Acknowledgments}

This work was conducted as part of the LEX AI platform development at legal.org.ua. Compute costs for all experiments were covered by an AWS Activate grant. We thank the creators of SIB-200, the ULP benchmark, and the EU Acts in Ukrainian corpus for making their datasets publicly available.

\paragraph{Data and code availability.} All datasets used are publicly available: SIB-200 \citep{adelani2024sib200} on HuggingFace, EU Acts in Ukrainian \citep{francophonicEU} on HuggingFace, ULP \citep{galeshchuk2024ulp} on HuggingFace. Experiment logs, prompts, and raw model outputs are available at \url{https://github.com/overthelex/SecondLayer}.

% ============================================================
% REFERENCES
% ============================================================
\begin{thebibliography}{20}
\providecommand{\natexlab}[1]{#1}

\bibitem[Adelani et~al.(2024)]{adelani2024sib200}
Adelani, D.~I., et~al.
\newblock SIB-200: A Simple, Inclusive, and Big Evaluation Dataset for Topic Classification in 200+ Languages and Dialects.
\newblock \emph{Proceedings of EACL}, pages 226--245, 2024.
\newblock \href{https://arxiv.org/abs/2309.07445}{arXiv:2309.07445}.

\bibitem[Ahia et~al.(2023)]{ahia2023all}
Ahia, O., Kumar, S., Gonen, H., Kasai, J., Mortensen, D.~R., Smith, N.~A., and Tsvetkov, Y.
\newblock Do All Languages Cost the Same? Tokenization in the Era of Commercial Language Models.
\newblock \emph{Proceedings of EMNLP}, pages 9904--9923, 2023.
\newblock \href{https://arxiv.org/abs/2305.13707}{arXiv:2305.13707}.

\bibitem[Brown et~al.(2020)]{brown2020language}
Brown, T.~B., et~al.
\newblock Language Models are Few-Shot Learners.
\newblock \emph{NeurIPS}, 2020.

\bibitem[Chaplynskyi(2023)]{chaplynskyi2023ukrbruk}
Chaplynskyi, D.
\newblock Ukrainian Brown Corpus and evaluation of multilingual models.
\newblock \emph{Proceedings of the Second Ukrainian NLP Workshop}, 2023.

\bibitem[Galeshchuk(2024)]{galeshchuk2024ulp}
Galeshchuk, S.
\newblock ULP: Ukrainian Language Proficiency Benchmark for LLMs.
\newblock \emph{HuggingFace Datasets}, 2024.

\bibitem[FrancophonIA(2022)]{francophonicEU}
FrancophonIA.
\newblock EU Acts in Ukrainian: Multilingual Legal Translation Corpus.
\newblock \emph{HuggingFace Datasets}, 2022.

\bibitem[Lai et~al.(2023)]{lai2023chatgpt}
Lai, V.~D., Ngo, N.~T., Veyseh, A.~P.~B., Man, H., Dernoncourt, F., Bui, T., and Nguyen, T.~H.
\newblock ChatGPT Beyond English: Towards a Comprehensive Evaluation of Large Language Models in Multilingual Learning.
\newblock \emph{arXiv preprint arXiv:2304.05613}, 2023.

\bibitem[Liu et~al.(2022)]{liu2022makes}
Liu, J., Shen, D., Zhang, Y., Dolan, B., Carin, L., and Chen, W.
\newblock What Makes Good In-Context Examples for GPT-3?
\newblock \emph{Proceedings of DeeLIO Workshop at ACL}, 2022.

\bibitem[Min et~al.(2022)]{min2022rethinking}
Min, S., Lyu, X., Holtzman, A., Artetxe, M., Lewis, M., Hajishirzi, H., and Zettlemoyer, L.
\newblock Rethinking the Role of Demonstrations: What Makes In-Context Learning Work?
\newblock \emph{Proceedings of EMNLP}, 2022.

\bibitem[Ovcharov(2026a)]{ovcharov2026tokenizer}
Ovcharov, V.
\newblock Tokenizer Fertility and Zero-Shot Performance of Foundation Models on Ukrainian Legal Text: A Comparative Study.
\newblock \href{https://arxiv.org/abs/2605.14890}{arXiv:2605.14890}, 2026.

\bibitem[Ovcharov(2026b)]{ovcharov2026fewshot}
Ovcharov, V.
\newblock Few-Shot Degradation Is Not What It Seems: Behavioral Evidence, Representation Analysis, and a Random-Text Control Across 12 Models, 2 Tasks, and 2 Architectures.
\newblock \emph{arXiv preprint}, 2026.

\bibitem[Petrov et~al.(2023)]{petrov2024language}
Petrov, A., La~Malfa, E., Torr, P., and Bibi, A.
\newblock Language Model Tokenizers Introduce Unfairness Between Languages.
\newblock \emph{NeurIPS}, 2023.
\newblock \href{https://arxiv.org/abs/2305.15425}{arXiv:2305.15425}.

\bibitem[Rust et~al.(2021)]{rust2021good}
Rust, P., Pfeiffer, J., Vuli\'{c}, I., Ruder, S., and Gurevych, I.
\newblock How Good is Your Tokenizer? On the Monolingual Performance of Multilingual Language Models.
\newblock \emph{Proceedings of ACL}, pages 3118--3135, 2021.

\bibitem[Abadji et~al.(2022)]{abadji2022towards}
Abadji, J., Ortiz~Suarez, P., Romary, L., and Sagot, B.
\newblock Towards a Cleaner Document-Oriented Multilingual Crawled Corpus.
\newblock \emph{Proceedings of LREC}, pages 4344--4355, 2022.

\bibitem[Nguyen et~al.(2024)]{nguyen2024culturax}
Nguyen, T., Nguyen, H., Lai, V.~D., Man, H., Ngo, N.~T., Dernoncourt, F., Rossi, R.~A., and Nguyen, T.~H.
\newblock {CulturaX}: A Cleaned, Enormous, and Multilingual Dataset for Large Language Models in 167~Languages.
\newblock \emph{Proceedings of ACL: System Demonstrations}, pages 228--237, 2024.
\newblock \href{https://arxiv.org/abs/2309.09400}{arXiv:2309.09400}.

\bibitem[Wei et~al.(2022)]{wei2022chain}
Wei, J., Wang, X., Schuurmans, D., Bosma, M., Ichter, B., Xia, F., Chi, E., Le, Q., and Zhou, D.
\newblock Chain-of-Thought Prompting Elicits Reasoning in Large Language Models.
\newblock \emph{NeurIPS}, 2022.

\bibitem[Conneau et~al.(2020)]{conneau2020unsupervised}
Conneau, A., Khandelwal, K., Goyal, N., Chaudhary, V., Wenzek, G., Guzm\'{a}n, F., Grave, E., Ott, M., Zettlemoyer, L., and Stoyanov, V.
\newblock Unsupervised Cross-lingual Representation Learning at Scale.
\newblock \emph{Proceedings of ACL}, pages 8440--8451, 2020.

\bibitem[Niklaus et~al.(2024)]{niklaus2024lextreme}
Niklaus, J., Chalkidis, I., and St\"urmer, M.
\newblock {LEXTREME}: A Multi-Lingual and Multi-Task Benchmark for the Legal Domain.
\newblock \emph{Findings of EACL}, pages 1721--1734, 2024.
\newblock \href{https://arxiv.org/abs/2301.13126}{arXiv:2301.13126}.

\bibitem[Zhao et~al.(2021)]{zhao2021calibrate}
Zhao, Z., Wallace, E., Feng, S., Klein, D., and Singh, S.
\newblock Calibrate Before Use: Improving Few-Shot Performance of Language Models.
\newblock \emph{Proceedings of ICML}, 2021.

\end{thebibliography}

% ============================================================
% APPENDIX
% ============================================================
\appendix

\section{EU Acts Fertility: Full Results}
\label{app:eu_fertility}

The full 25-language $\times$ 7-model fertility matrix is visualized in Figure~\ref{fig:eu_heatmap} (main text). Table~\ref{tab:eu_fertility} presents representative values for one model (Qwen3~32B); all seven models show the same language ranking with consistent offsets.

\end{document}
